% Recovered from the previously completed PDF; no new mathematical content.
% Prose is editable; formula boxes retain original symbol positions.
% Build from repository root: xelatex -no-shell-escape -halt-on-error -output-directory=... print/fudan/NAME.tex
\documentclass[11pt,a4paper]{article}
\usepackage[margin=23mm,headheight=14pt,headsep=8.5mm,footskip=10mm]{geometry}
\usepackage{amsmath,amssymb}
\DeclareMathAlphabet{\mathsf}{OT1}{cmss}{m}{n}
\usepackage[no-math]{fontspec}
\usepackage{xeCJK}
\setmainfont{Latin Modern Roman}
\setCJKmainfont{FandolSong-Regular.otf}[Path=assets/fonts/fandol/,BoldFont=FandolSong-Bold.otf,ItalicFont=FandolSong-Regular.otf,Script=CustomDefault]
\xeCJKsetup{PunctStyle=quanjiao,CheckSingle=true}
\usepackage{setspace}\setstretch{1.12}
\setlength{\parindent}{2em}\setlength{\parskip}{3pt}
\emergencystretch=6em
\hfuzz=9pt
\usepackage{titlesec}
\titleformat{\section}{\centering\normalfont\Large\bfseries}{}{0pt}{}
\titlespacing*{\section}{0pt}{1.8\baselineskip}{1.0\baselineskip}
\usepackage{fancyhdr}\pagestyle{fancy}\fancyhf{}
\fancyhead[L]{\small 数学手记·高等代数}\fancyhead[R]{\small 15A01–15A07}
\fancyfoot[C]{\thepage}\renewcommand{\headrulewidth}{0.4pt}
\usepackage[hidelinks,unicode]{hyperref}
\hypersetup{pdftitle={15A01–15A07：题目、解答与小结},pdfauthor={Vincent Arqain}}
\newcommand{\RB}[4]{\raisebox{0pt}[#2bp][#3bp]{\hbox to #1bp{#4\hss}}}
\newcommand{\RG}[5]{\hbox to 0pt{\kern#1bp\raisebox{#2bp}[0pt][0pt]{\fontsize{#3bp}{#3bp}\csname font#4\endcsname\char#5}\hss}}
\newcommand{\RT}[4]{\hbox to 0pt{\kern#1bp\raisebox{#2bp}[0pt][0pt]{\fontsize{#3bp}{#3bp}\selectfont#4}\hss}}
\newcommand{\RS}[4]{\hbox to 0pt{\kern#1bp\raisebox{#2bp}[0pt][0pt]{\fontsize{#3bp}{#3bp}\selectfont\ensuremath{#4}}\hss}}
\newcommand{\RR}[4]{\hbox to 0pt{\kern#1bp\raisebox{#2bp}[0pt][0pt]{\rule{#3bp}{#4bp}}\hss}}
\newcommand{\fontmi}{\usefont{OML}{cmm}{m}{it}}
\newcommand{\fontsy}{\usefont{OMS}{cmsy}{m}{n}}
\newcommand{\fontex}{\usefont{OMX}{cmex}{m}{n}}
\newcommand{\fontrm}{\usefont{OT1}{cmr}{m}{n}}
\newcommand{\fontbf}{\usefont{OT1}{cmr}{bx}{n}}
\newcommand{\fontam}{\usefont{U}{msa}{m}{n}}
\newcommand{\fontbb}{\usefont{U}{msb}{m}{n}}
\widowpenalty=10000\clubpenalty=10000\raggedbottom
\begin{document}
\begin{center}{\LARGE\bfseries 15A01–15A07\par}\vspace{8pt}复旦大学高等代数习题集·题目与解答\end{center}
% Original page 1.
复旦大学高等代数习题集题目与证明

题目选自谢启鸿《复旦大学高等代数习题集》（2023版），§1.4，第19–20页。证明与小结为本次整理。记号：\allowbreak{}\hskip.1em \RB{18.701}{7.655}{2.890}{%
\RG{0.000}{-0.004}{7.9701}{mi}{114}\RG{4.572}{-0.004}{7.9701}{rm}{40}\RG{8.284}{-0.004}{7.9701}{mi}{65}\RG{15.443}{-0.004}{7.9701}{rm}{41}%
}\hskip.1em\allowbreak{} 表示秩，\allowbreak{}\hskip.1em \RB{9.555}{7.505}{0.504}{%
\RG{0.000}{-0.004}{7.9701}{mi}{65}\RG{7.159}{3.461}{7.9701}{sy}{48}%
}\hskip.1em\allowbreak{} 表示转置，\allowbreak{}\hskip.1em \RB{10.804}{7.331}{0.504}{%
\RG{0.000}{-0.004}{7.9701}{mi}{65}\RG{7.159}{3.461}{7.9701}{sy}{3}%
}\hskip.1em\allowbreak{} 表示伴随矩阵；\allowbreak{}\hskip.1em \RB{67.151}{7.655}{2.890}{%
\RG{0.000}{-0.004}{7.9701}{bb}{78}\RG{9.545}{-0.004}{7.9701}{rm}{61}\RG{19.621}{-0.004}{7.9701}{sy}{102}\RG{24.394}{-0.004}{7.9701}{rm}{48}\RG{29.166}{-0.004}{7.9701}{mi}{59}\RG{33.409}{-0.004}{7.9701}{rm}{49}\RG{38.182}{-0.004}{7.9701}{mi}{59}\RG{42.424}{-0.004}{7.9701}{rm}{50}\RG{47.196}{-0.004}{7.9701}{mi}{59}\RG{51.435}{-0.004}{7.9701}{mi}{58}\RG{55.682}{-0.004}{7.9701}{mi}{58}\RG{59.930}{-0.004}{7.9701}{mi}{58}\RG{62.575}{-0.004}{7.9701}{sy}{103}%
}\hskip.1em\allowbreak{} 。

\section*{2015A01}
证明：第三类分块初等变换是若干个第三类初等变换的复合。特别地，第三类分块初等变换不改变行列式的值。

证明：设原矩阵的第\allowbreak{}\hskip.1em \RB{13.187}{7.664}{2.729}{%
\RG{0.000}{-0.004}{10.9091}{mi}{105}\RG{3.644}{-0.004}{10.9091}{mi}{59}\RG{8.422}{-0.004}{10.9091}{mi}{106}%
}\hskip.1em\allowbreak{} 分块行分别由第\allowbreak{}\hskip.1em \RB{44.655}{5.307}{3.546}{%
\RG{0.000}{-0.004}{10.9091}{mi}{117}\RG{6.079}{-1.640}{7.9701}{rm}{49}\RG{10.381}{-0.004}{10.9091}{mi}{59}\RG{15.159}{-0.004}{10.9091}{mi}{58}\RG{19.948}{-0.004}{10.9091}{mi}{58}\RG{24.737}{-0.004}{10.9091}{mi}{58}\RG{29.515}{-0.004}{10.9091}{mi}{59}\RG{34.305}{-0.004}{10.9091}{mi}{117}\RG{40.387}{-1.640}{7.9701}{mi}{112}%
}\hskip.1em\allowbreak{} 行和第\allowbreak{}\hskip.1em \RB{42.521}{5.307}{3.546}{%
\RG{0.000}{-0.004}{10.9091}{mi}{118}\RG{5.165}{-1.640}{7.9701}{rm}{49}\RG{9.466}{-0.004}{10.9091}{mi}{59}\RG{14.245}{-0.004}{10.9091}{mi}{58}\RG{19.034}{-0.004}{10.9091}{mi}{58}\RG{23.812}{-0.004}{10.9091}{mi}{58}\RG{28.601}{-0.004}{10.9091}{mi}{59}\RG{33.390}{-0.004}{10.9091}{mi}{118}\RG{38.557}{-1.640}{7.9701}{mi}{113}%
}\hskip.1em\allowbreak{} 行组成，记为\allowbreak{}\hskip.1em \RB{29.842}{7.947}{3.618}{%
\RG{0.000}{-0.004}{10.9091}{sy}{82}\RG{9.245}{-1.640}{7.9701}{mi}{105}\RG{12.319}{-0.004}{10.9091}{mi}{59}\RG{17.105}{-0.004}{10.9091}{sy}{82}\RG{26.350}{-1.640}{7.9701}{mi}{106}%
}\hskip.1em\allowbreak{} ；设\allowbreak{}\hskip.1em \RB{45.212}{8.678}{3.545}{%
\RG{0.000}{-0.004}{10.9091}{mi}{67}\RG{11.456}{-0.004}{10.9091}{rm}{61}\RG{22.791}{-0.004}{10.9091}{rm}{40}\RG{26.948}{-0.004}{10.9091}{mi}{99}\RG{31.545}{-1.639}{7.9701}{mi}{11}\RG{36.487}{-1.639}{7.9701}{mi}{12}\RG{41.635}{-0.004}{10.9091}{rm}{41}%
}\hskip.1em\allowbreak{} 是\allowbreak{}\hskip.1em \RB{24.005}{5.853}{2.620}{%
\RG{0.000}{-0.004}{10.9091}{mi}{112}\RG{7.785}{-0.004}{10.9091}{sy}{2}\RG{18.695}{-0.004}{10.9091}{mi}{113}%
}\hskip.1em\allowbreak{} 矩阵。第三类分块行变换\allowbreak{}\hskip.1em \RB{75.929}{8.176}{3.617}{%
\RG{0.000}{-0.004}{10.9091}{sy}{82}\RG{9.245}{-1.639}{7.9701}{mi}{105}\RG{15.350}{-0.004}{10.9091}{sy}{32}\RG{29.281}{-0.004}{10.9091}{sy}{82}\RG{38.535}{-1.639}{7.9701}{mi}{105}\RG{44.033}{-0.004}{10.9091}{rm}{43}\RG{54.765}{-0.004}{10.9091}{mi}{67}\RG{63.191}{-0.004}{10.9091}{sy}{82}\RG{72.436}{-1.639}{7.9701}{mi}{106}%
}\hskip.1em\allowbreak{} （\allowbreak{}\hskip.1em \RB{22.778}{8.307}{2.860}{%
\RG{0.000}{-0.004}{10.9091}{mi}{105}\RG{6.675}{-0.004}{10.9091}{sy}{54}\RG{6.675}{-0.004}{10.9091}{rm}{61}\RG{18.013}{-0.004}{10.9091}{mi}{106}%
}\hskip.1em\allowbreak{} ）要求

\[
\RB{196.297}{16.083}{12.965}{%
\RG{77.939}{12.388}{7.9701}{mi}{113}\RG{73.247}{9.067}{9.0569}{ex}{88}\RG{0.000}{0.001}{10.9091}{mi}{82}\RG{20.895}{0.001}{10.9091}{sy}{32}\RG{29.994}{0.001}{10.9091}{sy}{0}\RG{41.502}{0.001}{10.9091}{mi}{82}\RG{61.791}{0.001}{10.9091}{rm}{43}\RG{88.871}{0.001}{10.9091}{mi}{99}\RG{103.557}{0.001}{10.9091}{mi}{82}\RG{120.415}{0.001}{10.9091}{mi}{59}\RG{147.018}{0.001}{10.9091}{rm}{49}\RG{155.389}{0.001}{10.9091}{sy}{20}\RG{166.904}{0.001}{10.9091}{mi}{11}\RG{176.798}{0.001}{10.9091}{sy}{20}\RG{188.314}{0.001}{10.9091}{mi}{112}\RG{193.670}{0.001}{10.9091}{mi}{59}\RG{8.131}{-1.635}{7.9701}{mi}{117}\RG{49.632}{-1.635}{7.9701}{mi}{117}\RG{93.468}{-1.635}{7.9701}{mi}{11}\RG{98.410}{-1.635}{7.9701}{mi}{12}\RG{111.687}{-1.635}{7.9701}{mi}{118}\RG{12.588}{-2.542}{5.9776}{mi}{11}\RG{54.089}{-2.542}{5.9776}{mi}{11}\RG{115.436}{-2.914}{5.9776}{mi}{12}\RG{72.524}{-11.059}{7.9701}{mi}{12}\RG{77.215}{-11.059}{7.9701}{rm}{61}\RG{83.200}{-11.059}{7.9701}{rm}{49}%
}
\]
其中\allowbreak{}\hskip.1em \RB{12.833}{7.949}{2.211}{%
\RG{0.000}{-0.002}{10.9091}{mi}{82}\RG{8.131}{-1.638}{7.9701}{mi}{104}%
}\hskip.1em\allowbreak{} 表示原矩阵的第\allowbreak{}\hskip.1em \RB{6.336}{8.069}{0.611}{%
\RG{0.000}{-0.002}{10.9091}{mi}{104}%
}\hskip.1em\allowbreak{} 行。对每个\allowbreak{}\hskip.1em \RB{6.925}{5.309}{0.611}{%
\RG{0.000}{-0.002}{10.9091}{mi}{11}%
}\hskip.1em\allowbreak{} ，依次实施\allowbreak{}\hskip.1em \RB{97.212}{7.949}{4.471}{%
\RG{0.000}{-0.002}{10.9091}{mi}{82}\RG{8.130}{-1.638}{7.9701}{mi}{117}\RG{12.587}{-2.545}{5.9776}{mi}{11}\RG{21.133}{-0.002}{10.9091}{sy}{32}\RG{35.310}{-0.002}{10.9091}{mi}{82}\RG{43.440}{-1.638}{7.9701}{mi}{117}\RG{47.897}{-2.545}{5.9776}{mi}{11}\RG{55.694}{-0.002}{10.9091}{rm}{43}\RG{66.521}{-0.002}{10.9091}{mi}{99}\RG{71.118}{-1.638}{7.9701}{mi}{11}\RG{76.060}{-1.638}{7.9701}{mi}{12}\RG{81.208}{-0.002}{10.9091}{mi}{82}\RG{89.338}{-1.638}{7.9701}{mi}{118}\RG{93.087}{-2.916}{5.9776}{mi}{12}%
}\hskip.1em\allowbreak{} （\allowbreak{}\hskip.1em \RB{56.054}{8.189}{2.618}{%
\RG{0.000}{-0.002}{10.9091}{mi}{12}\RG{9.903}{-0.002}{10.9091}{rm}{61}\RG{21.477}{-0.002}{10.9091}{rm}{49}\RG{26.819}{-0.002}{10.9091}{mi}{59}\RG{31.597}{-0.002}{10.9091}{mi}{58}\RG{36.386}{-0.002}{10.9091}{mi}{58}\RG{41.175}{-0.002}{10.9091}{mi}{58}\RG{45.954}{-0.002}{10.9091}{mi}{59}\RG{50.743}{-0.002}{10.9091}{mi}{113}%
}\hskip.1em\allowbreak{} ，系数为零时略去）。两个分块行的指标互不相交，故来源行均未改变，最终每个目标行恰为上式，其余行不变，因而得到所需的分块行变换。对于分块列变换\allowbreak{}\hskip.1em \RB{66.883}{8.189}{3.616}{%
\RG{0.000}{-0.002}{10.9091}{sy}{67}\RG{5.745}{-1.638}{7.9701}{mi}{105}\RG{11.950}{-0.002}{10.9091}{sy}{32}\RG{25.990}{-0.002}{10.9091}{sy}{67}\RG{31.735}{-1.638}{7.9701}{mi}{105}\RG{37.275}{-0.002}{10.9091}{rm}{43}\RG{48.047}{-0.002}{10.9091}{sy}{67}\RG{53.791}{-1.638}{7.9701}{mi}{106}\RG{57.775}{-0.002}{10.9091}{mi}{68}%
}\hskip.1em\allowbreak{} ，转置后得到\allowbreak{}\hskip.1em \RB{69.197}{8.507}{4.677}{%
\RG{0.000}{-0.002}{10.9091}{sy}{67}\RG{5.744}{-2.699}{7.9701}{mi}{105}\RG{6.380}{3.956}{7.9701}{sy}{48}\RG{12.053}{-0.002}{10.9091}{sy}{32}\RG{26.093}{-0.002}{10.9091}{sy}{67}\RG{31.838}{-2.699}{7.9701}{mi}{105}\RG{32.475}{3.956}{7.9701}{sy}{48}\RG{37.481}{-0.002}{10.9091}{rm}{43}\RG{48.254}{-0.002}{10.9091}{mi}{68}\RG{57.419}{3.956}{7.9701}{sy}{48}\RG{59.960}{-0.002}{10.9091}{sy}{67}\RG{65.705}{-2.699}{7.9701}{mi}{106}\RG{66.341}{3.956}{7.9701}{sy}{48}%
}\hskip.1em\allowbreak{} 。应用行变换的结论再转置回来，即得列变换的结论。每次普通第三类初等变换都不改变方阵的行列式，因此其有限次复合也不改变行列式。\allowbreak{}\hskip.1em \RB{8.365}{7.993}{0.502}{%
\RG{0.000}{-0.002}{10.9091}{am}{3}%
}\hskip.1em\allowbreak{}

小结：把分块变换按系数拆成普通初等变换；来源行（列）始终不变，保证复合后的结果正好等于原分块变换。

\section*{2015A02}
设\allowbreak{}\hskip.1em \RB{56.699}{8.680}{3.616}{%
\RG{0.000}{-0.002}{10.9091}{mi}{65}\RG{11.038}{-0.002}{10.9091}{rm}{61}\RG{22.373}{-0.002}{10.9091}{rm}{40}\RG{26.530}{-0.002}{10.9091}{mi}{97}\RG{32.137}{-1.638}{7.9701}{mi}{105}\RG{34.760}{-1.638}{7.9701}{mi}{106}\RG{38.743}{-0.002}{10.9091}{rm}{40}\RG{42.896}{-0.002}{10.9091}{mi}{120}\RG{48.966}{-0.002}{10.9091}{rm}{41}\RG{53.123}{-0.002}{10.9091}{rm}{41}%
}\hskip.1em\allowbreak{} 是\allowbreak{}\hskip.1em \RB{6.555}{5.309}{0.611}{%
\RG{0.000}{-0.002}{10.9091}{mi}{110}%
}\hskip.1em\allowbreak{} （\allowbreak{}\hskip.1em \RB{26.222}{7.753}{1.996}{%
\RG{0.000}{-0.002}{10.9091}{mi}{110}\RG{9.406}{-0.002}{10.9091}{sy}{21}\RG{20.922}{-0.002}{10.9091}{rm}{50}%
}\hskip.1em\allowbreak{} ）阶方阵，其中每个元素\allowbreak{}\hskip.1em \RB{26.012}{8.680}{3.616}{%
\RG{0.000}{-0.002}{10.9091}{mi}{97}\RG{5.607}{-1.638}{7.9701}{mi}{105}\RG{8.230}{-1.638}{7.9701}{mi}{106}\RG{12.212}{-0.002}{10.9091}{rm}{40}\RG{16.365}{-0.002}{10.9091}{mi}{120}\RG{22.435}{-0.002}{10.9091}{rm}{41}%
}\hskip.1em\allowbreak{} 都是关于未定元\allowbreak{}\hskip.1em \RB{6.107}{5.309}{0.611}{%
\RG{0.000}{-0.002}{10.9091}{mi}{120}%
}\hskip.1em\allowbreak{} 的多项式。若\allowbreak{}\hskip.1em \RB{5.965}{8.069}{0.611}{%
\RG{0.000}{-0.002}{10.9091}{mi}{107}%
}\hskip.1em\allowbreak{} 是正整数，满足\allowbreak{}\hskip.1em \RB{10.432}{9.486}{0.611}{%
\RG{0.000}{-0.002}{10.9091}{mi}{120}\RG{6.070}{3.957}{7.9701}{mi}{107}%
}\hskip.1em\allowbreak{} 整除\allowbreak{}\hskip.1em \RB{8.202}{8.287}{0.502}{%
\RG{0.000}{-0.002}{10.9091}{mi}{65}%
}\hskip.1em\allowbreak{} 的所有代数余子式\allowbreak{}\hskip.1em \RB{14.123}{8.287}{3.616}{%
\RG{0.000}{-0.002}{10.9091}{mi}{65}\RG{8.007}{-1.638}{7.9701}{mi}{105}\RG{10.630}{-1.638}{7.9701}{mi}{106}%
}\hskip.1em\allowbreak{} ，证明：\allowbreak{}\hskip.1em \RB{19.995}{9.486}{0.611}{%
\RG{0.000}{-0.002}{10.9091}{mi}{120}\RG{6.070}{3.957}{7.9701}{mi}{107}\RG{10.271}{3.957}{7.9701}{rm}{43}\RG{16.256}{3.957}{7.9701}{rm}{49}%
}\hskip.1em\allowbreak{} 整除\allowbreak{}\hskip.1em \RB{13.272}{8.680}{3.229}{%
\RG{0.000}{-0.002}{10.9091}{sy}{106}\RG{3.030}{-0.002}{10.9091}{mi}{65}\RG{11.037}{-0.002}{10.9091}{sy}{106}%
}\hskip.1em\allowbreak{} 。

证明：记\allowbreak{}\hskip.1em \RB{76.374}{8.680}{3.240}{%
\RG{0.000}{-0.002}{10.9091}{mi}{68}\RG{9.165}{-0.002}{10.9091}{rm}{40}\RG{13.320}{-0.002}{10.9091}{mi}{120}\RG{19.390}{-0.002}{10.9091}{rm}{41}\RG{26.568}{-0.002}{10.9091}{rm}{61}\RG{37.914}{-0.002}{10.9091}{rm}{100}\RG{43.848}{-0.002}{10.9091}{rm}{101}\RG{48.594}{-0.002}{10.9091}{rm}{116}\RG{54.566}{-0.002}{10.9091}{mi}{65}\RG{62.574}{-0.002}{10.9091}{rm}{40}\RG{66.727}{-0.002}{10.9091}{mi}{120}\RG{72.797}{-0.002}{10.9091}{rm}{41}%
}\hskip.1em\allowbreak{} 。若\allowbreak{}\hskip.1em \RB{43.300}{8.680}{3.240}{%
\RG{0.000}{-0.002}{10.9091}{mi}{68}\RG{9.165}{-0.002}{10.9091}{rm}{40}\RG{13.319}{-0.002}{10.9091}{mi}{120}\RG{19.389}{-0.002}{10.9091}{rm}{41}\RG{26.567}{-0.002}{10.9091}{rm}{61}\RG{37.913}{-0.002}{10.9091}{rm}{48}%
}\hskip.1em\allowbreak{} ，结论成立。以下设\allowbreak{}\hskip.1em \RB{43.295}{8.680}{3.240}{%
\RG{0.000}{-0.002}{10.9091}{mi}{68}\RG{9.165}{-0.002}{10.9091}{rm}{40}\RG{13.319}{-0.002}{10.9091}{mi}{120}\RG{19.390}{-0.002}{10.9091}{rm}{41}\RG{26.574}{-0.002}{10.9091}{sy}{54}\RG{26.574}{-0.002}{10.9091}{rm}{61}\RG{37.908}{-0.002}{10.9091}{rm}{48}%
}\hskip.1em\allowbreak{} 。由\allowbreak{}\hskip.1em \RB{67.874}{8.680}{3.240}{%
\RG{0.000}{-0.002}{10.9091}{mi}{65}\RG{8.007}{-0.002}{10.9091}{mi}{65}\RG{16.015}{3.957}{7.9701}{sy}{3}\RG{23.347}{-0.002}{10.9091}{rm}{61}\RG{34.685}{-0.002}{10.9091}{mi}{68}\RG{43.851}{-0.002}{10.9091}{rm}{40}\RG{48.005}{-0.002}{10.9091}{mi}{120}\RG{54.075}{-0.002}{10.9091}{rm}{41}\RG{58.229}{-0.002}{10.9091}{mi}{73}\RG{62.940}{-1.638}{7.9701}{mi}{110}%
}\hskip.1em\allowbreak{}取行列式，并在多项式环中约去非零多项式\allowbreak{}\hskip.1em \RB{22.965}{8.680}{3.240}{%
\RG{0.000}{-0.002}{10.9091}{mi}{68}\RG{9.165}{-0.002}{10.9091}{rm}{40}\RG{13.319}{-0.002}{10.9091}{mi}{120}\RG{19.389}{-0.002}{10.9091}{rm}{41}%
}\hskip.1em\allowbreak{} ，得\allowbreak{}\hskip.1em \RB{84.546}{9.276}{3.240}{%
\RG{0.000}{-0.002}{10.9091}{rm}{100}\RG{5.935}{-0.002}{10.9091}{rm}{101}\RG{10.680}{-0.002}{10.9091}{rm}{116}\RG{16.654}{-0.002}{10.9091}{mi}{65}\RG{24.662}{3.957}{7.9701}{sy}{3}\RG{33.630}{-0.002}{10.9091}{rm}{61}\RG{46.605}{-0.002}{10.9091}{mi}{68}\RG{55.771}{-0.002}{10.9091}{rm}{40}\RG{59.925}{-0.002}{10.9091}{mi}{120}\RG{65.995}{-0.002}{10.9091}{rm}{41}\RG{70.148}{3.957}{7.9701}{mi}{110}\RG{74.819}{3.957}{7.9701}{sy}{0}\RG{80.807}{3.957}{7.9701}{rm}{49}%
}\hskip.1em\allowbreak{} 。又由假设可写\allowbreak{}\hskip.1em \RB{59.868}{9.486}{3.240}{%
\RG{0.000}{-0.002}{10.9091}{mi}{65}\RG{8.007}{3.957}{7.9701}{sy}{3}\RG{15.340}{-0.002}{10.9091}{rm}{61}\RG{26.678}{-0.002}{10.9091}{mi}{120}\RG{32.748}{3.957}{7.9701}{mi}{107}\RG{37.403}{-0.002}{10.9091}{mi}{66}\RG{46.068}{-0.002}{10.9091}{rm}{40}\RG{50.222}{-0.002}{10.9091}{mi}{120}\RG{56.292}{-0.002}{10.9091}{rm}{41}%
}\hskip.1em\allowbreak{} ，其中\allowbreak{}\hskip.1em \RB{22.465}{8.680}{3.240}{%
\RG{0.000}{-0.002}{10.9091}{mi}{66}\RG{8.665}{-0.002}{10.9091}{rm}{40}\RG{12.819}{-0.002}{10.9091}{mi}{120}\RG{18.889}{-0.002}{10.9091}{rm}{41}%
}\hskip.1em\allowbreak{} 为多项式矩阵，故

\[
\RB{155.624}{10.035}{3.236}{%
\RG{23.544}{4.506}{7.9701}{mi}{110}\RG{28.215}{4.506}{7.9701}{sy}{0}\RG{34.202}{4.506}{7.9701}{rm}{49}\RG{77.535}{4.506}{7.9701}{sy}{3}\RG{102.275}{4.506}{7.9701}{mi}{107}\RG{106.471}{4.506}{7.9701}{mi}{110}\RG{0.000}{0.003}{10.9091}{mi}{68}\RG{9.165}{0.003}{10.9091}{rm}{40}\RG{13.319}{0.003}{10.9091}{mi}{120}\RG{19.389}{0.003}{10.9091}{rm}{41}\RG{41.535}{0.003}{10.9091}{rm}{61}\RG{52.869}{0.003}{10.9091}{rm}{100}\RG{58.804}{0.003}{10.9091}{rm}{101}\RG{63.549}{0.003}{10.9091}{rm}{116}\RG{69.527}{0.003}{10.9091}{mi}{65}\RG{84.866}{0.003}{10.9091}{rm}{61}\RG{96.205}{0.003}{10.9091}{mi}{120}\RG{113.419}{0.003}{10.9091}{rm}{100}\RG{119.354}{0.003}{10.9091}{rm}{101}\RG{124.099}{0.003}{10.9091}{rm}{116}\RG{130.073}{0.003}{10.9091}{mi}{66}\RG{138.738}{0.003}{10.9091}{rm}{40}\RG{142.892}{0.003}{10.9091}{mi}{120}\RG{148.962}{0.003}{10.9091}{rm}{41}\RG{153.116}{0.003}{10.9091}{mi}{58}%
}
\]
设\allowbreak{}\hskip.1em \RB{71.194}{8.680}{3.240}{%
\RG{0.000}{-0.002}{10.9091}{mi}{68}\RG{9.165}{-0.002}{10.9091}{rm}{40}\RG{13.320}{-0.002}{10.9091}{mi}{120}\RG{19.390}{-0.002}{10.9091}{rm}{41}\RG{26.568}{-0.002}{10.9091}{rm}{61}\RG{37.912}{-0.002}{10.9091}{mi}{120}\RG{43.983}{3.956}{7.9701}{mi}{109}\RG{51.245}{-0.002}{10.9091}{mi}{104}\RG{57.394}{-0.002}{10.9091}{rm}{40}\RG{61.547}{-0.002}{10.9091}{mi}{120}\RG{67.617}{-0.002}{10.9091}{rm}{41}%
}\hskip.1em\allowbreak{} ，其中\allowbreak{}\hskip.1em \RB{39.549}{8.680}{3.240}{%
\RG{0.000}{-0.002}{10.9091}{mi}{104}\RG{6.148}{-0.002}{10.9091}{rm}{40}\RG{10.305}{-0.002}{10.9091}{rm}{48}\RG{15.650}{-0.002}{10.9091}{rm}{41}\RG{22.827}{-0.002}{10.9091}{sy}{54}\RG{22.827}{-0.002}{10.9091}{rm}{61}\RG{34.162}{-0.002}{10.9091}{rm}{48}%
}\hskip.1em\allowbreak{} 。比较上式两边因子\allowbreak{}\hskip.1em \RB{6.107}{5.309}{0.611}{%
\RG{0.000}{-0.002}{10.9091}{mi}{120}%
}\hskip.1em\allowbreak{} 的重数，得

\[
\RB{125.701}{8.684}{3.236}{%
\RG{0.000}{0.002}{10.9091}{rm}{40}\RG{4.155}{0.002}{10.9091}{mi}{110}\RG{12.955}{0.002}{10.9091}{sy}{0}\RG{23.864}{0.002}{10.9091}{rm}{49}\RG{29.209}{0.002}{10.9091}{rm}{41}\RG{33.359}{0.002}{10.9091}{mi}{109}\RG{45.732}{0.002}{10.9091}{sy}{21}\RG{57.247}{0.002}{10.9091}{mi}{107}\RG{63.160}{0.002}{10.9091}{mi}{110}\RG{72.575}{0.002}{10.9091}{mi}{62}\RG{83.920}{0.002}{10.9091}{mi}{107}\RG{89.835}{0.002}{10.9091}{rm}{40}\RG{93.989}{0.002}{10.9091}{mi}{110}\RG{102.790}{0.002}{10.9091}{sy}{0}\RG{113.699}{0.002}{10.9091}{rm}{49}\RG{119.045}{0.002}{10.9091}{rm}{41}\RG{123.194}{0.002}{10.9091}{mi}{58}%
}
\]
因\allowbreak{}\hskip.1em \RB{26.222}{7.752}{1.997}{%
\RG{0.000}{-0.003}{10.9091}{mi}{110}\RG{9.406}{-0.003}{10.9091}{sy}{21}\RG{20.922}{-0.003}{10.9091}{rm}{50}%
}\hskip.1em\allowbreak{} ，故\allowbreak{}\hskip.1em \RB{29.682}{8.068}{0.808}{%
\RG{0.000}{-0.003}{10.9091}{mi}{109}\RG{12.371}{-0.003}{10.9091}{mi}{62}\RG{23.716}{-0.003}{10.9091}{mi}{107}%
}\hskip.1em\allowbreak{} ；又\allowbreak{}\hskip.1em \RB{9.522}{5.308}{0.612}{%
\RG{0.000}{-0.003}{10.9091}{mi}{109}%
}\hskip.1em\allowbreak{} 是整数，所以\allowbreak{}\hskip.1em \RB{47.942}{8.068}{1.997}{%
\RG{0.000}{-0.003}{10.9091}{mi}{109}\RG{12.374}{-0.003}{10.9091}{sy}{21}\RG{23.889}{-0.003}{10.9091}{mi}{107}\RG{32.235}{-0.003}{10.9091}{rm}{43}\RG{42.969}{-0.003}{10.9091}{rm}{49}%
}\hskip.1em\allowbreak{} ，即\allowbreak{}\hskip.1em \RB{52.618}{9.485}{3.241}{%
\RG{0.000}{-0.003}{10.9091}{mi}{120}\RG{6.070}{3.956}{7.9701}{mi}{107}\RG{10.272}{3.956}{7.9701}{rm}{43}\RG{16.257}{3.956}{7.9701}{rm}{49}\RG{23.591}{-0.003}{10.9091}{sy}{106}\RG{29.652}{-0.003}{10.9091}{mi}{68}\RG{38.817}{-0.003}{10.9091}{rm}{40}\RG{42.972}{-0.003}{10.9091}{mi}{120}\RG{49.042}{-0.003}{10.9091}{rm}{41}%
}\hskip.1em\allowbreak{} 。\allowbreak{}\hskip.1em \RB{8.365}{7.992}{0.503}{%
\RG{0.000}{-0.003}{10.9091}{am}{3}%
}\hskip.1em\allowbreak{}

小结：用伴随矩阵集中所有代数余子式的整除条件，再比较因子\allowbreak{}\hskip.1em \RB{6.107}{5.308}{0.612}{%
\RG{0.000}{-0.003}{10.9091}{mi}{120}%
}\hskip.1em\allowbreak{} 的重数。最后利用重数是整数，得到至少\allowbreak{}\hskip.1em \RB{24.053}{8.068}{1.343}{%
\RG{0.000}{-0.003}{10.9091}{mi}{107}\RG{8.345}{-0.003}{10.9091}{rm}{43}\RG{19.080}{-0.003}{10.9091}{rm}{49}%
}\hskip.1em\allowbreak{} 次整除。

% Original page 2.
\section*{2015A03}
设

\[
\RB{205.267}{10.134}{56.058}{%
\RG{25.842}{9.281}{9.0569}{ex}{48}\RG{192.897}{9.281}{9.0569}{ex}{49}\RG{53.988}{3.955}{7.9701}{rm}{50}\RG{48.381}{-0.004}{10.9091}{mi}{97}\RG{82.373}{-0.004}{10.9091}{mi}{97}\RG{92.282}{-0.004}{10.9091}{mi}{97}\RG{104.615}{-0.004}{10.9091}{rm}{43}\RG{115.350}{-0.004}{10.9091}{rm}{49}\RG{130.156}{-0.004}{10.9091}{sy}{1}\RG{135.000}{-0.004}{10.9091}{sy}{1}\RG{139.855}{-0.004}{10.9091}{sy}{1}\RG{153.759}{-0.004}{10.9091}{mi}{97}\RG{163.668}{-0.004}{10.9091}{mi}{97}\RG{176.824}{-0.004}{10.9091}{rm}{43}\RG{187.558}{-0.004}{10.9091}{rm}{49}\RG{87.980}{-1.640}{7.9701}{rm}{49}\RG{97.889}{-1.640}{7.9701}{rm}{50}\RG{159.366}{-1.640}{7.9701}{rm}{49}\RG{169.275}{-1.640}{7.9701}{mi}{110}\RG{53.988}{-2.701}{7.9701}{rm}{49}\RG{25.842}{-6.660}{9.0569}{ex}{66}\RG{192.897}{-6.660}{9.0569}{ex}{67}\RG{25.842}{-12.094}{9.0569}{ex}{66}\RG{192.897}{-12.095}{9.0569}{ex}{67}\RG{102.184}{-12.982}{7.9701}{rm}{50}\RG{34.177}{-16.940}{10.9091}{mi}{97}\RG{44.087}{-16.940}{10.9091}{mi}{97}\RG{56.421}{-16.940}{10.9091}{rm}{43}\RG{67.155}{-16.940}{10.9091}{rm}{49}\RG{96.576}{-16.940}{10.9091}{mi}{97}\RG{130.156}{-16.940}{10.9091}{sy}{1}\RG{135.000}{-16.940}{10.9091}{sy}{1}\RG{139.855}{-16.940}{10.9091}{sy}{1}\RG{153.759}{-16.940}{10.9091}{mi}{97}\RG{163.668}{-16.940}{10.9091}{mi}{97}\RG{176.824}{-16.940}{10.9091}{rm}{43}\RG{187.558}{-16.940}{10.9091}{rm}{49}\RG{25.842}{-17.528}{9.0569}{ex}{66}\RG{192.897}{-17.528}{9.0569}{ex}{67}\RG{39.785}{-18.576}{7.9701}{rm}{50}\RG{49.695}{-18.576}{7.9701}{rm}{49}\RG{159.366}{-18.576}{7.9701}{rm}{50}\RG{169.275}{-18.576}{7.9701}{mi}{110}\RG{102.184}{-19.637}{7.9701}{rm}{50}\RG{25.842}{-22.963}{9.0569}{ex}{66}\RG{192.897}{-22.963}{9.0569}{ex}{67}\RG{0.000}{-25.690}{10.9091}{mi}{77}\RG{14.504}{-25.690}{10.9091}{rm}{61}\RG{202.640}{-25.690}{10.9091}{mi}{59}\RG{25.842}{-28.396}{9.0569}{ex}{66}\RG{192.897}{-28.396}{9.0569}{ex}{67}\RT{51.852}{-28.514}{10.9091}{.}\RT{100.047}{-28.514}{10.9091}{.}\RT{171.845}{-28.514}{10.9091}{.}\RT{131.766}{-29.419}{10.9091}{.}\RT{51.852}{-32.136}{10.9091}{.}\RT{100.047}{-32.136}{10.9091}{.}\RT{135.945}{-32.136}{10.9091}{.}\RT{171.845}{-32.136}{10.9091}{.}\RG{25.842}{-33.831}{9.0569}{ex}{66}\RG{192.897}{-33.831}{9.0569}{ex}{67}\RT{140.125}{-34.854}{10.9091}{.}\RT{51.852}{-35.759}{10.9091}{.}\RT{100.047}{-35.759}{10.9091}{.}\RT{171.845}{-35.759}{10.9091}{.}\RG{25.842}{-39.627}{9.0569}{ex}{64}\RG{192.897}{-39.627}{9.0569}{ex}{65}\RG{173.570}{-48.736}{7.9701}{rm}{50}\RG{33.766}{-52.695}{10.9091}{mi}{97}\RG{44.498}{-52.695}{10.9091}{mi}{97}\RG{56.832}{-52.695}{10.9091}{rm}{43}\RG{67.566}{-52.695}{10.9091}{rm}{49}\RG{81.962}{-52.695}{10.9091}{mi}{97}\RG{92.693}{-52.695}{10.9091}{mi}{97}\RG{105.026}{-52.695}{10.9091}{rm}{43}\RG{115.761}{-52.695}{10.9091}{rm}{49}\RG{130.156}{-52.695}{10.9091}{sy}{1}\RG{135.000}{-52.695}{10.9091}{sy}{1}\RG{139.855}{-52.695}{10.9091}{sy}{1}\RG{167.963}{-52.695}{10.9091}{mi}{97}\RG{39.374}{-54.332}{7.9701}{mi}{110}\RG{50.105}{-54.332}{7.9701}{rm}{49}\RG{87.569}{-54.332}{7.9701}{mi}{110}\RG{98.300}{-54.332}{7.9701}{rm}{50}\RG{173.570}{-55.393}{7.9701}{mi}{110}%
}
\]
证明：\allowbreak{}\hskip.1em \RB{64.119}{8.684}{3.236}{%
\RG{0.000}{0.002}{10.9091}{mi}{114}\RG{5.111}{0.002}{10.9091}{rm}{40}\RG{9.265}{0.002}{10.9091}{mi}{77}\RG{20.738}{0.002}{10.9091}{rm}{41}\RG{27.922}{0.002}{10.9091}{sy}{21}\RG{39.437}{0.002}{10.9091}{mi}{110}\RG{48.237}{0.002}{10.9091}{sy}{0}\RG{59.146}{0.002}{10.9091}{rm}{49}%
}\hskip.1em\allowbreak{} 。

证明：\allowbreak{}\hskip.1em \RB{25.714}{7.756}{0.607}{%
\RG{0.000}{0.002}{10.9091}{mi}{110}\RG{9.406}{0.002}{10.9091}{rm}{61}\RG{20.741}{0.002}{10.9091}{rm}{49}%
}\hskip.1em\allowbreak{} 时结论显然。以下设\allowbreak{}\hskip.1em \RB{26.222}{7.756}{1.993}{%
\RG{0.000}{0.002}{10.9091}{mi}{110}\RG{9.406}{0.002}{10.9091}{sy}{21}\RG{20.922}{0.002}{10.9091}{rm}{50}%
}\hskip.1em\allowbreak{} ，令\allowbreak{}\hskip.1em \RB{75.336}{8.684}{3.236}{%
\RG{0.000}{0.002}{10.9091}{mi}{97}\RG{8.637}{0.002}{10.9091}{rm}{61}\RG{19.972}{0.002}{10.9091}{rm}{40}\RG{24.130}{0.002}{10.9091}{mi}{97}\RG{29.737}{-1.635}{7.9701}{rm}{49}\RG{34.039}{0.002}{10.9091}{mi}{59}\RG{38.817}{0.002}{10.9091}{mi}{58}\RG{43.606}{0.002}{10.9091}{mi}{58}\RG{48.395}{0.002}{10.9091}{mi}{58}\RG{53.174}{0.002}{10.9091}{mi}{59}\RG{57.963}{0.002}{10.9091}{mi}{97}\RG{63.573}{-1.635}{7.9701}{mi}{110}\RG{68.697}{0.002}{10.9091}{rm}{41}\RG{72.851}{3.961}{7.9701}{sy}{48}%
}\hskip.1em\allowbreak{} 、\allowbreak{}\hskip.1em \RB{64.721}{8.684}{3.236}{%
\RG{0.000}{0.002}{10.9091}{mi}{101}\RG{7.981}{0.002}{10.9091}{rm}{61}\RG{19.315}{0.002}{10.9091}{rm}{40}\RG{23.472}{0.002}{10.9091}{rm}{49}\RG{28.814}{0.002}{10.9091}{mi}{59}\RG{33.592}{0.002}{10.9091}{mi}{58}\RG{38.381}{0.002}{10.9091}{mi}{58}\RG{43.170}{0.002}{10.9091}{mi}{58}\RG{47.948}{0.002}{10.9091}{mi}{59}\RG{52.740}{0.002}{10.9091}{rm}{49}\RG{58.085}{0.002}{10.9091}{rm}{41}\RG{62.235}{3.961}{7.9701}{sy}{48}%
}\hskip.1em\allowbreak{} 、\allowbreak{}\hskip.1em \RB{33.670}{8.511}{0.727}{%
\RG{0.000}{0.002}{10.9091}{mi}{74}\RG{9.945}{0.002}{10.9091}{rm}{61}\RG{21.283}{0.002}{10.9091}{mi}{101}\RG{26.235}{0.002}{10.9091}{mi}{101}\RG{31.185}{3.961}{7.9701}{sy}{48}%
}\hskip.1em\allowbreak{} ，则\allowbreak{}\hskip.1em \RB{82.647}{8.511}{2.207}{%
\RG{0.000}{0.002}{10.9091}{mi}{77}\RG{14.504}{0.002}{10.9091}{rm}{61}\RG{25.843}{0.002}{10.9091}{mi}{97}\RG{31.450}{0.002}{10.9091}{mi}{97}\RG{37.057}{3.961}{7.9701}{sy}{48}\RG{42.022}{0.002}{10.9091}{rm}{43}\RG{52.755}{0.002}{10.9091}{mi}{74}\RG{62.093}{0.002}{10.9091}{sy}{0}\RG{73.002}{0.002}{10.9091}{mi}{73}\RG{77.713}{-1.635}{7.9701}{mi}{110}%
}\hskip.1em\allowbreak{} 。由\allowbreak{}\hskip.1em \RB{39.219}{9.279}{0.727}{%
\RG{0.000}{0.002}{10.9091}{mi}{74}\RG{6.914}{3.961}{7.9701}{rm}{50}\RG{14.246}{0.002}{10.9091}{rm}{61}\RG{25.585}{0.002}{10.9091}{mi}{110}\RG{31.966}{0.002}{10.9091}{mi}{74}%
}\hskip.1em\allowbreak{} ，直接展开可得

\[
\RB{137.680}{15.135}{8.632}{%
\RG{40.208}{13.234}{9.0569}{ex}{18}\RG{104.182}{13.234}{9.0569}{ex}{19}\RG{57.817}{7.380}{10.9091}{rm}{49}\RG{0.000}{0.000}{10.9091}{rm}{40}\RG{4.154}{0.000}{10.9091}{mi}{74}\RG{13.492}{0.000}{10.9091}{sy}{0}\RG{24.401}{0.000}{10.9091}{mi}{73}\RG{34.236}{0.000}{10.9091}{rm}{41}\RG{74.099}{0.000}{10.9091}{mi}{74}\RG{83.437}{0.000}{10.9091}{sy}{0}\RG{94.346}{0.000}{10.9091}{mi}{73}\RG{113.879}{-0.000}{10.9091}{rm}{61}\RG{125.217}{-0.000}{10.9091}{mi}{73}\RG{135.053}{-0.000}{10.9091}{mi}{59}\RG{29.113}{-1.636}{7.9701}{mi}{110}\RG{99.058}{-1.636}{7.9701}{mi}{110}\RG{129.929}{-1.636}{7.9701}{mi}{110}\RG{47.962}{-7.483}{10.9091}{mi}{110}\RG{56.763}{-7.484}{10.9091}{sy}{0}\RG{67.672}{-7.483}{10.9091}{rm}{49}\RR{47.962}{2.727}{25.051}{0.362}%
}
\]
故\allowbreak{}\hskip.1em \RB{29.893}{7.954}{2.206}{%
\RG{0.000}{0.003}{10.9091}{mi}{74}\RG{9.338}{0.003}{10.9091}{sy}{0}\RG{20.247}{0.003}{10.9091}{mi}{73}\RG{24.959}{-1.634}{7.9701}{mi}{110}%
}\hskip.1em\allowbreak{} 是非异阵。由\allowbreak{}\hskip.1em \RB{82.958}{8.512}{2.206}{%
\RG{0.000}{0.003}{10.9091}{mi}{74}\RG{9.338}{0.003}{10.9091}{sy}{0}\RG{20.247}{0.003}{10.9091}{mi}{73}\RG{24.958}{-1.634}{7.9701}{mi}{110}\RG{33.113}{0.003}{10.9091}{rm}{61}\RG{44.451}{0.003}{10.9091}{mi}{77}\RG{58.349}{0.003}{10.9091}{sy}{0}\RG{69.258}{0.003}{10.9091}{mi}{97}\RG{74.865}{0.003}{10.9091}{mi}{97}\RG{80.473}{3.962}{7.9701}{sy}{48}%
}\hskip.1em\allowbreak{} 及秩的不等式，

\[
\RB{204.575}{9.053}{3.240}{%
\RG{137.318}{4.503}{7.9701}{sy}{48}\RG{0.000}{-0.002}{10.9091}{mi}{110}\RG{9.406}{-0.002}{10.9091}{rm}{61}\RG{20.745}{-0.002}{10.9091}{mi}{114}\RG{25.855}{-0.002}{10.9091}{rm}{40}\RG{30.009}{-0.002}{10.9091}{mi}{74}\RG{39.347}{-0.002}{10.9091}{sy}{0}\RG{50.256}{-0.002}{10.9091}{mi}{73}\RG{60.092}{-0.002}{10.9091}{rm}{41}\RG{67.275}{-0.002}{10.9091}{sy}{20}\RG{78.791}{-0.002}{10.9091}{mi}{114}\RG{83.901}{-0.002}{10.9091}{rm}{40}\RG{88.055}{-0.002}{10.9091}{mi}{77}\RG{99.528}{-0.002}{10.9091}{rm}{41}\RG{106.106}{-0.002}{10.9091}{rm}{43}\RG{116.839}{-0.002}{10.9091}{mi}{114}\RG{121.949}{-0.002}{10.9091}{rm}{40}\RG{126.104}{-0.002}{10.9091}{mi}{97}\RG{131.711}{-0.002}{10.9091}{mi}{97}\RG{139.859}{-0.002}{10.9091}{rm}{41}\RG{147.043}{-0.002}{10.9091}{sy}{20}\RG{158.558}{-0.002}{10.9091}{mi}{114}\RG{163.669}{-0.002}{10.9091}{rm}{40}\RG{167.823}{-0.002}{10.9091}{mi}{77}\RG{179.296}{-0.002}{10.9091}{rm}{41}\RG{185.875}{-0.002}{10.9091}{rm}{43}\RG{196.609}{-0.002}{10.9091}{rm}{49}\RG{201.947}{-0.002}{10.9091}{mi}{59}\RG{54.967}{-1.638}{7.9701}{mi}{110}%
}
\]
从而\allowbreak{}\hskip.1em \RB{64.120}{8.685}{3.235}{%
\RG{0.000}{0.003}{10.9091}{mi}{114}\RG{5.111}{0.003}{10.9091}{rm}{40}\RG{9.265}{0.003}{10.9091}{mi}{77}\RG{20.738}{0.003}{10.9091}{rm}{41}\RG{27.922}{0.003}{10.9091}{sy}{21}\RG{39.437}{0.003}{10.9091}{mi}{110}\RG{48.238}{0.003}{10.9091}{sy}{0}\RG{59.147}{0.003}{10.9091}{rm}{49}%
}\hskip.1em\allowbreak{} 。\allowbreak{}\hskip.1em \RB{8.365}{7.997}{0.500}{%
\RG{0.000}{0.003}{10.9091}{am}{3}%
}\hskip.1em\allowbreak{}

小结：将\allowbreak{}\hskip.1em \RB{11.605}{7.954}{0.500}{%
\RG{0.000}{0.003}{10.9091}{mi}{77}%
}\hskip.1em\allowbreak{} 写成“非异阵\allowbreak{}\hskip.1em \RB{8.213}{6.786}{1.337}{%
\RG{0.000}{0.003}{10.9091}{rm}{43}%
}\hskip.1em\allowbreak{} 秩至多为\allowbreak{}\hskip.1em \RB{4.973}{7.757}{0.500}{%
\RG{0.000}{0.003}{10.9091}{rm}{49}%
}\hskip.1em\allowbreak{} 的矩阵”，即可得到所需的秩下界，不必直接计算\allowbreak{}\hskip.1em \RB{11.605}{7.954}{0.500}{%
\RG{0.000}{0.003}{10.9091}{mi}{77}%
}\hskip.1em\allowbreak{} 的子式。

\section*{2015A04}
设\allowbreak{}\hskip.1em \RB{8.202}{8.292}{0.500}{%
\RG{0.000}{0.003}{10.9091}{mi}{65}%
}\hskip.1em\allowbreak{} 是\allowbreak{}\hskip.1em \RB{29.232}{5.859}{0.606}{%
\RG{0.000}{0.003}{10.9091}{mi}{109}\RG{11.768}{0.003}{10.9091}{sy}{2}\RG{22.677}{0.003}{10.9091}{mi}{110}%
}\hskip.1em\allowbreak{} 实矩阵，试用秩的子式判别法和Cauchy–Binet公式证明：

\[
\RB{116.619}{9.052}{3.240}{%
\RG{17.272}{4.502}{7.9701}{sy}{48}\RG{71.623}{4.502}{7.9701}{sy}{48}\RG{0.000}{-0.002}{10.9091}{mi}{114}\RG{5.111}{-0.002}{10.9091}{rm}{40}\RG{9.265}{-0.002}{10.9091}{mi}{65}\RG{19.814}{-0.002}{10.9091}{mi}{65}\RG{27.821}{-0.002}{10.9091}{rm}{41}\RG{34.999}{-0.002}{10.9091}{rm}{61}\RG{46.344}{-0.002}{10.9091}{mi}{114}\RG{51.454}{-0.002}{10.9091}{rm}{40}\RG{55.607}{-0.002}{10.9091}{mi}{65}\RG{63.615}{-0.002}{10.9091}{mi}{65}\RG{74.164}{-0.002}{10.9091}{rm}{41}\RG{81.342}{-0.002}{10.9091}{rm}{61}\RG{92.686}{-0.002}{10.9091}{mi}{114}\RG{97.796}{-0.002}{10.9091}{rm}{40}\RG{101.951}{-0.002}{10.9091}{mi}{65}\RG{109.958}{-0.002}{10.9091}{rm}{41}\RG{114.112}{-0.002}{10.9091}{mi}{58}%
}
\]
证明：设\allowbreak{}\hskip.1em \RB{40.942}{8.684}{3.236}{%
\RG{0.000}{0.002}{10.9091}{mi}{114}\RG{5.111}{0.002}{10.9091}{rm}{40}\RG{9.265}{0.002}{10.9091}{mi}{65}\RG{17.272}{0.002}{10.9091}{rm}{41}\RG{24.450}{0.002}{10.9091}{rm}{61}\RG{35.795}{0.002}{10.9091}{mi}{114}%
}\hskip.1em\allowbreak{} 。若\allowbreak{}\hskip.1em \RB{24.862}{7.756}{0.727}{%
\RG{0.000}{0.002}{10.9091}{mi}{114}\RG{8.140}{0.002}{10.9091}{rm}{61}\RG{19.475}{0.002}{10.9091}{rm}{48}%
}\hskip.1em\allowbreak{} ，则\allowbreak{}\hskip.1em \RB{30.796}{8.291}{0.727}{%
\RG{0.000}{0.002}{10.9091}{mi}{65}\RG{11.037}{0.002}{10.9091}{rm}{61}\RG{22.376}{0.002}{10.9091}{mi}{79}%
}\hskip.1em\allowbreak{} ，结论成立。以下设\allowbreak{}\hskip.1em \RB{24.628}{7.756}{1.993}{%
\RG{0.000}{0.002}{10.9091}{mi}{114}\RG{8.140}{0.002}{10.9091}{sy}{21}\RG{19.655}{0.002}{10.9091}{rm}{49}%
}\hskip.1em\allowbreak{} 。用\allowbreak{}\hskip.1em \RB{30.497}{8.684}{3.225}{%
\RG{0.000}{0.002}{10.9091}{mi}{65}\RG{8.007}{0.002}{10.9091}{rm}{91}\RG{11.038}{0.002}{10.9091}{mi}{73}\RG{16.602}{0.002}{10.9091}{mi}{59}\RG{21.391}{0.002}{10.9091}{mi}{74}\RG{28.306}{0.002}{10.9091}{rm}{93}%
}\hskip.1em\allowbreak{} 表示按递增指标集\allowbreak{}\hskip.1em \RB{17.605}{7.953}{2.604}{%
\RG{0.000}{0.002}{10.9091}{mi}{73}\RG{5.564}{0.002}{10.9091}{mi}{59}\RG{10.353}{0.002}{10.9091}{mi}{74}%
}\hskip.1em\allowbreak{} 取出的子矩阵。由秩的子式判别法，存在\allowbreak{}\hskip.1em \RB{69.337}{8.684}{3.225}{%
\RG{0.000}{0.002}{10.9091}{mi}{73}\RG{8.599}{0.002}{10.9091}{sy}{18}\RG{20.108}{0.002}{10.9091}{sy}{102}\RG{25.568}{0.002}{10.9091}{rm}{49}\RG{30.909}{0.002}{10.9091}{mi}{59}\RG{35.687}{0.002}{10.9091}{mi}{58}\RG{40.476}{0.002}{10.9091}{mi}{58}\RG{45.265}{0.002}{10.9091}{mi}{58}\RG{50.044}{0.002}{10.9091}{mi}{59}\RG{54.833}{0.002}{10.9091}{mi}{109}\RG{64.179}{0.002}{10.9091}{sy}{103}%
}\hskip.1em\allowbreak{} 、\allowbreak{}\hskip.1em \RB{67.716}{8.685}{3.225}{%
\RG{0.000}{0.002}{10.9091}{mi}{74}\RG{9.945}{0.003}{10.9091}{sy}{18}\RG{21.454}{0.003}{10.9091}{sy}{102}\RG{26.915}{0.002}{10.9091}{rm}{49}\RG{32.255}{0.003}{10.9091}{mi}{59}\RG{37.034}{0.003}{10.9091}{mi}{58}\RG{41.823}{0.003}{10.9091}{mi}{58}\RG{46.601}{0.003}{10.9091}{mi}{58}\RG{51.390}{0.003}{10.9091}{mi}{59}\RG{56.179}{0.003}{10.9091}{mi}{110}\RG{62.558}{0.002}{10.9091}{sy}{103}%
}\hskip.1em\allowbreak{} ，均含\allowbreak{}\hskip.1em \RB{5.147}{5.313}{0.607}{%
\RG{0.000}{0.002}{10.9091}{mi}{114}%
}\hskip.1em\allowbreak{} 个指标，使得\allowbreak{}\hskip.1em \RB{67.679}{8.684}{3.225}{%
\RG{0.000}{0.002}{10.9091}{rm}{100}\RG{5.935}{0.002}{10.9091}{rm}{101}\RG{10.680}{0.002}{10.9091}{rm}{116}\RG{16.654}{0.003}{10.9091}{mi}{65}\RG{24.662}{0.002}{10.9091}{rm}{91}\RG{27.692}{0.003}{10.9091}{mi}{73}\RG{33.255}{0.003}{10.9091}{mi}{59}\RG{38.045}{0.003}{10.9091}{mi}{74}\RG{44.960}{0.002}{10.9091}{rm}{93}\RG{50.957}{0.003}{10.9091}{sy}{54}\RG{50.957}{0.002}{10.9091}{rm}{61}\RG{62.292}{0.002}{10.9091}{rm}{48}%
}\hskip.1em\allowbreak{} 。由Cauchy–Binet公式，

\[
\RB{290.060}{11.383}{61.835}{%
\RG{14.835}{10.521}{9.0569}{ex}{16}\RG{69.029}{10.521}{9.0569}{ex}{17}\RG{101.259}{9.072}{9.0569}{ex}{88}\RG{40.414}{4.508}{7.9701}{sy}{48}\RG{205.126}{4.508}{7.9701}{sy}{48}\RG{0.000}{0.005}{10.9091}{rm}{100}\RG{5.935}{0.005}{10.9091}{rm}{101}\RG{10.680}{0.005}{10.9091}{rm}{116}\RG{20.245}{0.005}{10.9091}{rm}{40}\RG{24.399}{0.005}{10.9091}{mi}{65}\RG{32.406}{0.005}{10.9091}{mi}{65}\RG{42.955}{0.005}{10.9091}{rm}{41}\RG{47.111}{0.005}{10.9091}{rm}{91}\RG{50.139}{0.005}{10.9091}{mi}{73}\RG{55.703}{0.005}{10.9091}{mi}{59}\RG{60.492}{0.005}{10.9091}{mi}{73}\RG{66.062}{0.005}{10.9091}{rm}{93}\RG{77.468}{0.005}{10.9091}{rm}{61}\RG{128.611}{0.005}{10.9091}{rm}{100}\RG{134.545}{0.005}{10.9091}{rm}{101}\RG{139.291}{0.005}{10.9091}{rm}{116}\RG{145.265}{0.005}{10.9091}{mi}{65}\RG{153.273}{0.005}{10.9091}{rm}{91}\RG{156.303}{0.005}{10.9091}{mi}{73}\RG{161.866}{0.005}{10.9091}{mi}{59}\RG{166.655}{0.005}{10.9091}{mi}{83}\RG{173.861}{0.005}{10.9091}{rm}{93}\RG{180.461}{0.005}{10.9091}{rm}{100}\RG{186.395}{0.005}{10.9091}{rm}{101}\RG{191.141}{0.005}{10.9091}{rm}{116}\RG{197.119}{0.005}{10.9091}{mi}{65}\RG{207.667}{0.005}{10.9091}{rm}{91}\RG{210.697}{0.005}{10.9091}{mi}{83}\RG{217.297}{0.005}{10.9091}{mi}{59}\RG{222.086}{0.005}{10.9091}{mi}{73}\RG{227.663}{0.005}{10.9091}{rm}{93}\RG{88.806}{-11.458}{7.9701}{mi}{83}\RG{93.894}{-11.458}{7.9701}{sy}{18}\RG{99.878}{-11.458}{7.9701}{sy}{102}\RG{103.731}{-11.458}{7.9701}{rm}{49}\RG{107.580}{-11.458}{7.9701}{mi}{59}\RG{109.717}{-11.458}{7.9701}{mi}{58}\RG{111.855}{-11.458}{7.9701}{mi}{58}\RG{113.992}{-11.458}{7.9701}{mi}{58}\RG{116.130}{-11.458}{7.9701}{mi}{59}\RG{118.267}{-11.458}{7.9701}{mi}{110}\RG{122.944}{-11.458}{7.9701}{sy}{103}\RG{98.280}{-19.790}{7.9701}{sy}{106}\RG{100.418}{-19.790}{7.9701}{mi}{83}\RG{105.505}{-19.790}{7.9701}{sy}{106}\RG{107.644}{-19.790}{7.9701}{rm}{61}\RG{113.632}{-19.790}{7.9701}{mi}{114}\RG{126.793}{-29.213}{9.0569}{ex}{16}\RG{180.419}{-29.213}{9.0569}{ex}{17}\RG{204.675}{-29.213}{9.0569}{ex}{16}\RG{258.012}{-29.213}{9.0569}{ex}{17}\RG{101.258}{-30.662}{9.0569}{ex}{88}\RG{185.828}{-31.718}{7.9701}{rm}{50}\RG{263.421}{-31.718}{7.9701}{rm}{50}\RG{77.468}{-39.729}{10.9091}{rm}{61}\RG{132.202}{-39.729}{10.9091}{rm}{100}\RG{138.136}{-39.729}{10.9091}{rm}{101}\RG{142.882}{-39.729}{10.9091}{rm}{116}\RG{148.855}{-39.729}{10.9091}{mi}{65}\RG{156.863}{-39.729}{10.9091}{rm}{91}\RG{159.894}{-39.729}{10.9091}{mi}{73}\RG{165.457}{-39.729}{10.9091}{mi}{59}\RG{170.246}{-39.729}{10.9091}{mi}{83}\RG{177.452}{-39.729}{10.9091}{rm}{93}\RG{193.161}{-39.729}{10.9091}{sy}{21}\RG{210.085}{-39.729}{10.9091}{rm}{100}\RG{216.019}{-39.729}{10.9091}{rm}{101}\RG{220.765}{-39.729}{10.9091}{rm}{116}\RG{226.739}{-39.729}{10.9091}{mi}{65}\RG{234.746}{-39.729}{10.9091}{rm}{91}\RG{237.776}{-39.729}{10.9091}{mi}{73}\RG{243.340}{-39.729}{10.9091}{mi}{59}\RG{248.129}{-39.729}{10.9091}{mi}{74}\RG{255.045}{-39.729}{10.9091}{rm}{93}\RG{270.754}{-39.729}{10.9091}{mi}{62}\RG{282.092}{-39.729}{10.9091}{rm}{48}\RG{287.433}{-39.729}{10.9091}{mi}{59}\RG{88.806}{-51.191}{7.9701}{mi}{83}\RG{93.894}{-51.191}{7.9701}{sy}{18}\RG{99.878}{-51.191}{7.9701}{sy}{102}\RG{103.731}{-51.191}{7.9701}{rm}{49}\RG{107.580}{-51.191}{7.9701}{mi}{59}\RG{109.717}{-51.191}{7.9701}{mi}{58}\RG{111.855}{-51.191}{7.9701}{mi}{58}\RG{113.992}{-51.191}{7.9701}{mi}{58}\RG{116.130}{-51.191}{7.9701}{mi}{59}\RG{118.267}{-51.191}{7.9701}{mi}{110}\RG{122.944}{-51.191}{7.9701}{sy}{103}\RG{98.280}{-59.524}{7.9701}{sy}{106}\RG{100.418}{-59.524}{7.9701}{mi}{83}\RG{105.505}{-59.524}{7.9701}{sy}{106}\RG{107.644}{-59.524}{7.9701}{rm}{61}\RG{113.632}{-59.524}{7.9701}{mi}{114}%
}
\]
其中用到\allowbreak{}\hskip.1em \RB{81.929}{8.680}{3.229}{%
\RG{0.000}{-0.002}{10.9091}{mi}{65}\RG{8.007}{3.957}{7.9701}{sy}{48}\RG{10.548}{-0.002}{10.9091}{rm}{91}\RG{13.579}{-0.002}{10.9091}{mi}{83}\RG{20.179}{-0.002}{10.9091}{mi}{59}\RG{24.968}{-0.002}{10.9091}{mi}{73}\RG{30.544}{-0.002}{10.9091}{rm}{93}\RG{36.533}{-0.002}{10.9091}{rm}{61}\RG{47.880}{-0.002}{10.9091}{mi}{65}\RG{55.887}{-0.002}{10.9091}{rm}{91}\RG{58.917}{-0.002}{10.9091}{mi}{73}\RG{64.481}{-0.002}{10.9091}{mi}{59}\RG{69.270}{-0.002}{10.9091}{mi}{83}\RG{76.476}{-0.002}{10.9091}{rm}{93}\RG{79.444}{3.957}{7.9701}{sy}{48}%
}\hskip.1em\allowbreak{} 及转置不改变行列式。故\allowbreak{}\hskip.1em \RB{18.501}{8.507}{0.502}{%
\RG{0.000}{-0.002}{10.9091}{mi}{65}\RG{8.007}{-0.002}{10.9091}{mi}{65}\RG{16.015}{3.957}{7.9701}{sy}{48}%
}\hskip.1em\allowbreak{} 有非零的\allowbreak{}\hskip.1em \RB{5.147}{5.309}{0.611}{%
\RG{0.000}{-0.002}{10.9091}{mi}{114}%
}\hskip.1em\allowbreak{} 阶子式，从而\allowbreak{}\hskip.1em \RB{43.075}{8.680}{3.240}{%
\RG{0.000}{-0.002}{10.9091}{mi}{114}\RG{5.110}{-0.002}{10.9091}{rm}{40}\RG{9.265}{-0.002}{10.9091}{mi}{65}\RG{17.272}{-0.002}{10.9091}{mi}{65}\RG{25.279}{3.957}{7.9701}{sy}{48}\RG{27.820}{-0.002}{10.9091}{rm}{41}\RG{35.005}{-0.002}{10.9091}{sy}{21}%
}\hskip.1em\allowbreak{}\allowbreak{}\hskip.1em \RB{5.147}{5.309}{0.611}{%
\RG{0.000}{-0.002}{10.9091}{mi}{114}%
}\hskip.1em\allowbreak{} 。另一方面，\allowbreak{}\hskip.1em \RB{87.461}{8.680}{3.240}{%
\RG{0.000}{-0.002}{10.9091}{mi}{114}\RG{5.110}{-0.002}{10.9091}{rm}{40}\RG{9.264}{-0.002}{10.9091}{mi}{65}\RG{17.271}{-0.002}{10.9091}{mi}{65}\RG{25.279}{3.957}{7.9701}{sy}{48}\RG{27.820}{-0.002}{10.9091}{rm}{41}\RG{35.005}{-0.002}{10.9091}{sy}{20}\RG{46.519}{-0.002}{10.9091}{mi}{114}\RG{51.630}{-0.002}{10.9091}{rm}{40}\RG{55.784}{-0.002}{10.9091}{mi}{65}\RG{63.791}{-0.002}{10.9091}{rm}{41}\RG{70.969}{-0.002}{10.9091}{rm}{61}\RG{82.314}{-0.002}{10.9091}{mi}{114}%
}\hskip.1em\allowbreak{} ，所以\allowbreak{}\hskip.1em \RB{67.191}{8.680}{3.240}{%
\RG{0.000}{-0.002}{10.9091}{mi}{114}\RG{5.110}{-0.002}{10.9091}{rm}{40}\RG{9.264}{-0.002}{10.9091}{mi}{65}\RG{17.271}{-0.002}{10.9091}{mi}{65}\RG{25.279}{3.957}{7.9701}{sy}{48}\RG{27.820}{-0.002}{10.9091}{rm}{41}\RG{34.998}{-0.002}{10.9091}{rm}{61}\RG{46.343}{-0.002}{10.9091}{mi}{114}\RG{51.453}{-0.002}{10.9091}{rm}{40}\RG{55.607}{-0.002}{10.9091}{mi}{65}\RG{63.615}{-0.002}{10.9091}{rm}{41}%
}\hskip.1em\allowbreak{} 。将这一结论应用于实矩阵\allowbreak{}\hskip.1em \RB{10.493}{8.507}{0.502}{%
\RG{0.000}{-0.002}{10.9091}{mi}{65}\RG{8.007}{3.957}{7.9701}{sy}{48}%
}\hskip.1em\allowbreak{} ，便有\allowbreak{}\hskip.1em \RB{105.527}{8.680}{3.240}{%
\RG{0.000}{-0.002}{10.9091}{mi}{114}\RG{5.111}{-0.002}{10.9091}{rm}{40}\RG{9.265}{-0.002}{10.9091}{mi}{65}\RG{17.272}{3.957}{7.9701}{sy}{48}\RG{19.814}{-0.002}{10.9091}{mi}{65}\RG{27.821}{-0.002}{10.9091}{rm}{41}\RG{34.999}{-0.002}{10.9091}{rm}{61}\RG{46.344}{-0.002}{10.9091}{mi}{114}\RG{51.454}{-0.002}{10.9091}{rm}{40}\RG{55.607}{-0.002}{10.9091}{mi}{65}\RG{63.615}{3.957}{7.9701}{sy}{48}\RG{66.156}{-0.002}{10.9091}{rm}{41}\RG{73.335}{-0.002}{10.9091}{rm}{61}\RG{84.679}{-0.002}{10.9091}{mi}{114}\RG{89.789}{-0.002}{10.9091}{rm}{40}\RG{93.944}{-0.002}{10.9091}{mi}{65}\RG{101.951}{-0.002}{10.9091}{rm}{41}%
}\hskip.1em\allowbreak{} ，于是所需等式全部成立。\allowbreak{}\hskip.1em \RB{8.365}{7.993}{0.502}{%
\RG{0.000}{-0.002}{10.9091}{am}{3}%
}\hskip.1em\allowbreak{}

小结：Cauchy–Binet公式把选定的主子式写成若干实数平方之和，其中至少一项非零。“实矩阵”保证各项不会相消。

% Original page 3.
\section*{2015A05}
设\allowbreak{}\hskip.1em \RB{21.391}{8.289}{2.605}{%
\RG{0.000}{0.000}{10.9091}{mi}{65}\RG{8.007}{0.000}{10.9091}{mi}{59}\RG{12.785}{0.000}{10.9091}{mi}{66}%
}\hskip.1em\allowbreak{} 都是\allowbreak{}\hskip.1em \RB{6.555}{5.311}{0.609}{%
\RG{0.000}{0.000}{10.9091}{mi}{110}%
}\hskip.1em\allowbreak{} 阶方阵，形式上约定\allowbreak{}\hskip.1em \RB{36.323}{9.277}{2.209}{%
\RG{0.000}{0.000}{10.9091}{mi}{65}\RG{8.007}{3.959}{7.9701}{rm}{48}\RG{15.340}{0.000}{10.9091}{rm}{61}\RG{26.678}{0.000}{10.9091}{mi}{73}\RG{31.389}{-1.636}{7.9701}{mi}{110}%
}\hskip.1em\allowbreak{} 。

（1）若\allowbreak{}\hskip.1em \RB{5.965}{8.068}{0.612}{%
\RG{0.000}{-0.003}{10.9091}{mi}{107}%
}\hskip.1em\allowbreak{} 是非负整数，使得\allowbreak{}\hskip.1em \RB{75.788}{9.485}{3.241}{%
\RG{0.000}{-0.003}{10.9091}{mi}{114}\RG{5.110}{-0.003}{10.9091}{rm}{40}\RG{9.265}{-0.003}{10.9091}{mi}{65}\RG{17.272}{3.956}{7.9701}{mi}{107}\RG{21.926}{-0.003}{10.9091}{rm}{41}\RG{29.105}{-0.003}{10.9091}{rm}{61}\RG{40.449}{-0.003}{10.9091}{mi}{114}\RG{45.559}{-0.003}{10.9091}{rm}{40}\RG{49.714}{-0.003}{10.9091}{mi}{65}\RG{57.721}{3.956}{7.9701}{mi}{107}\RG{61.922}{3.956}{7.9701}{rm}{43}\RG{67.907}{3.956}{7.9701}{rm}{49}\RG{72.212}{-0.003}{10.9091}{rm}{41}%
}\hskip.1em\allowbreak{} ，证明：对任意的\allowbreak{}\hskip.1em \RB{24.155}{8.068}{1.997}{%
\RG{0.000}{-0.003}{10.9091}{mi}{105}\RG{6.675}{-0.003}{10.9091}{sy}{21}\RG{18.190}{-0.003}{10.9091}{mi}{107}%
}\hskip.1em\allowbreak{} ，有\allowbreak{}\hskip.1em \RB{64.371}{9.485}{3.241}{%
\RG{0.000}{-0.003}{10.9091}{mi}{114}\RG{5.111}{-0.003}{10.9091}{rm}{40}\RG{9.265}{-0.003}{10.9091}{mi}{65}\RG{17.272}{3.956}{7.9701}{mi}{105}\RG{20.346}{-0.003}{10.9091}{rm}{41}\RG{27.525}{-0.003}{10.9091}{rm}{61}\RG{38.869}{-0.003}{10.9091}{mi}{114}\RG{43.979}{-0.003}{10.9091}{rm}{40}\RG{48.133}{-0.003}{10.9091}{mi}{65}\RG{56.141}{3.956}{7.9701}{mi}{107}\RG{60.795}{-0.003}{10.9091}{rm}{41}%
}\hskip.1em\allowbreak{} 。（2）记

\[
\RB{179.635}{10.030}{3.241}{%
\RG{112.459}{4.502}{7.9701}{mi}{107}\RG{152.908}{4.502}{7.9701}{mi}{107}\RG{157.109}{4.502}{7.9701}{rm}{43}\RG{163.094}{4.502}{7.9701}{rm}{49}\RG{0.000}{-0.003}{10.9091}{mi}{115}\RG{5.031}{-0.003}{10.9091}{rm}{40}\RG{9.185}{-0.003}{10.9091}{mi}{65}\RG{17.193}{-0.003}{10.9091}{rm}{41}\RG{24.371}{-0.003}{10.9091}{rm}{61}\RG{35.716}{-0.003}{10.9091}{rm}{109}\RG{44.618}{-0.003}{10.9091}{rm}{105}\RG{47.585}{-0.003}{10.9091}{rm}{110}\RG{53.571}{-0.003}{10.9091}{sy}{102}\RG{59.025}{-0.003}{10.9091}{mi}{107}\RG{67.977}{-0.003}{10.9091}{sy}{50}\RG{78.281}{-0.003}{10.9091}{bb}{78}\RG{89.190}{-0.003}{10.9091}{rm}{58}\RG{95.187}{-0.003}{10.9091}{mi}{114}\RG{100.297}{-0.003}{10.9091}{rm}{40}\RG{104.452}{-0.003}{10.9091}{mi}{65}\RG{117.114}{-0.003}{10.9091}{rm}{41}\RG{124.292}{-0.003}{10.9091}{rm}{61}\RG{135.636}{-0.003}{10.9091}{mi}{114}\RG{140.746}{-0.003}{10.9091}{rm}{40}\RG{144.900}{-0.003}{10.9091}{mi}{65}\RG{167.399}{-0.003}{10.9091}{rm}{41}\RG{171.553}{-0.003}{10.9091}{sy}{103}\RG{177.007}{-0.003}{10.9091}{mi}{59}%
}
\]
称为\allowbreak{}\hskip.1em \RB{8.202}{8.286}{0.503}{%
\RG{0.000}{-0.003}{10.9091}{mi}{65}%
}\hskip.1em\allowbreak{} 的稳定指数；最终稳定的秩记为\allowbreak{}\hskip.1em \RB{28.703}{8.679}{3.241}{%
\RG{0.000}{-0.003}{10.9091}{mi}{114}\RG{4.814}{-1.639}{7.9701}{sy}{49}\RG{12.965}{-0.003}{10.9091}{rm}{40}\RG{17.118}{-0.003}{10.9091}{mi}{65}\RG{25.126}{-0.003}{10.9091}{rm}{41}%
}\hskip.1em\allowbreak{} ，即\allowbreak{}\hskip.1em \RB{67.572}{9.253}{3.241}{%
\RG{0.000}{-0.003}{10.9091}{mi}{114}\RG{4.814}{-1.639}{7.9701}{sy}{49}\RG{12.965}{-0.003}{10.9091}{rm}{40}\RG{17.119}{-0.003}{10.9091}{mi}{65}\RG{25.126}{-0.003}{10.9091}{rm}{41}\RG{32.305}{-0.003}{10.9091}{rm}{61}\RG{43.649}{-0.003}{10.9091}{mi}{114}\RG{48.760}{-0.003}{10.9091}{rm}{40}\RG{52.914}{-0.003}{10.9091}{mi}{65}\RG{60.921}{3.956}{7.9701}{mi}{105}\RG{63.995}{-0.003}{10.9091}{rm}{41}%
}\hskip.1em\allowbreak{} （\allowbreak{}\hskip.1em \RB{38.958}{8.679}{3.241}{%
\RG{0.000}{-0.003}{10.9091}{mi}{105}\RG{6.674}{-0.003}{10.9091}{sy}{21}\RG{18.189}{-0.003}{10.9091}{mi}{115}\RG{23.221}{-0.003}{10.9091}{rm}{40}\RG{27.375}{-0.003}{10.9091}{mi}{65}\RG{35.382}{-0.003}{10.9091}{rm}{41}%
}\hskip.1em\allowbreak{} ）。证明：\allowbreak{}\hskip.1em \RB{20.769}{8.679}{3.241}{%
\RG{0.000}{-0.003}{10.9091}{mi}{115}\RG{5.032}{-0.003}{10.9091}{rm}{40}\RG{9.185}{-0.003}{10.9091}{mi}{65}\RG{17.193}{-0.003}{10.9091}{rm}{41}%
}\hskip.1em\allowbreak{} 必存在，且\allowbreak{}\hskip.1em \RB{62.333}{8.679}{3.241}{%
\RG{0.000}{-0.003}{10.9091}{rm}{48}\RG{8.371}{-0.003}{10.9091}{sy}{20}\RG{19.886}{-0.003}{10.9091}{mi}{115}\RG{24.917}{-0.003}{10.9091}{rm}{40}\RG{29.072}{-0.003}{10.9091}{mi}{65}\RG{37.079}{-0.003}{10.9091}{rm}{41}\RG{44.264}{-0.003}{10.9091}{sy}{20}\RG{55.778}{-0.003}{10.9091}{mi}{110}%
}\hskip.1em\allowbreak{} 。（3）证明：\allowbreak{}\hskip.1em \RB{89.683}{8.679}{3.241}{%
\RG{0.000}{-0.003}{10.9091}{mi}{114}\RG{4.814}{-1.639}{7.9701}{sy}{49}\RG{12.965}{-0.003}{10.9091}{rm}{40}\RG{17.119}{-0.003}{10.9091}{mi}{65}\RG{25.126}{-0.003}{10.9091}{mi}{66}\RG{33.792}{-0.003}{10.9091}{rm}{41}\RG{40.970}{-0.003}{10.9091}{rm}{61}\RG{52.315}{-0.003}{10.9091}{mi}{114}\RG{57.128}{-1.639}{7.9701}{sy}{49}\RG{65.280}{-0.003}{10.9091}{rm}{40}\RG{69.434}{-0.003}{10.9091}{mi}{66}\RG{78.095}{-0.003}{10.9091}{mi}{65}\RG{86.106}{-0.003}{10.9091}{rm}{41}%
}\hskip.1em\allowbreak{} 。（4）证明：\allowbreak{}\hskip.1em \RB{98.936}{8.679}{3.241}{%
\RG{0.000}{-0.003}{10.9091}{sy}{106}\RG{3.031}{-0.003}{10.9091}{mi}{115}\RG{8.062}{-0.003}{10.9091}{rm}{40}\RG{12.216}{-0.003}{10.9091}{mi}{65}\RG{20.224}{-0.003}{10.9091}{mi}{66}\RG{28.889}{-0.003}{10.9091}{rm}{41}\RG{35.467}{-0.003}{10.9091}{sy}{0}\RG{46.376}{-0.003}{10.9091}{mi}{115}\RG{51.407}{-0.003}{10.9091}{rm}{40}\RG{55.562}{-0.003}{10.9091}{mi}{66}\RG{64.224}{-0.003}{10.9091}{mi}{65}\RG{72.235}{-0.003}{10.9091}{rm}{41}\RG{76.388}{-0.003}{10.9091}{sy}{106}\RG{82.443}{-0.003}{10.9091}{sy}{20}\RG{93.964}{-0.003}{10.9091}{rm}{49}%
}\hskip.1em\allowbreak{} ，并举例说明可取到\allowbreak{}\hskip.1em \RB{21.391}{8.286}{2.608}{%
\RG{0.000}{-0.003}{10.9091}{mi}{65}\RG{8.007}{-0.003}{10.9091}{mi}{59}\RG{12.785}{-0.003}{10.9091}{mi}{66}%
}\hskip.1em\allowbreak{} ，使得\allowbreak{}\hskip.1em \RB{98.756}{8.679}{3.241}{%
\RG{0.000}{-0.003}{10.9091}{sy}{106}\RG{3.031}{-0.003}{10.9091}{mi}{115}\RG{8.062}{-0.003}{10.9091}{rm}{40}\RG{12.216}{-0.003}{10.9091}{mi}{65}\RG{20.224}{-0.003}{10.9091}{mi}{66}\RG{28.889}{-0.003}{10.9091}{rm}{41}\RG{35.467}{-0.003}{10.9091}{sy}{0}\RG{46.376}{-0.003}{10.9091}{mi}{115}\RG{51.407}{-0.003}{10.9091}{rm}{40}\RG{55.562}{-0.003}{10.9091}{mi}{66}\RG{64.224}{-0.003}{10.9091}{mi}{65}\RG{72.235}{-0.003}{10.9091}{rm}{41}\RG{76.388}{-0.003}{10.9091}{sy}{106}\RG{82.449}{-0.003}{10.9091}{rm}{61}\RG{93.784}{-0.003}{10.9091}{rm}{49}%
}\hskip.1em\allowbreak{} 。

证明：（1）记\allowbreak{}\hskip.1em \RB{22.675}{7.955}{0.500}{%
\RG{0.000}{0.004}{10.9091}{rm}{73}\RG{3.851}{0.004}{10.9091}{rm}{109}\RG{14.626}{0.004}{10.9091}{mi}{84}%
}\hskip.1em\allowbreak{} 为矩阵\allowbreak{}\hskip.1em \RB{8.049}{7.900}{0.500}{%
\RG{0.000}{0.004}{10.9091}{mi}{84}%
}\hskip.1em\allowbreak{} 的列空间。由于\allowbreak{}\hskip.1em \RB{82.693}{9.491}{1.991}{%
\RG{0.000}{0.004}{10.9091}{rm}{73}\RG{3.851}{0.004}{10.9091}{rm}{109}\RG{14.625}{0.004}{10.9091}{mi}{65}\RG{22.633}{3.963}{7.9701}{mi}{107}\RG{26.835}{3.963}{7.9701}{rm}{43}\RG{32.819}{3.963}{7.9701}{rm}{49}\RG{42.168}{0.004}{10.9091}{sy}{18}\RG{55.698}{0.004}{10.9091}{rm}{73}\RG{59.549}{0.004}{10.9091}{rm}{109}\RG{70.324}{0.004}{10.9091}{mi}{65}\RG{78.331}{3.963}{7.9701}{mi}{107}%
}\hskip.1em\allowbreak{} ，且两者维数相同，故\allowbreak{}\hskip.1em \RB{78.488}{9.491}{0.500}{%
\RG{0.000}{0.004}{10.9091}{rm}{73}\RG{3.851}{0.004}{10.9091}{rm}{109}\RG{14.625}{0.004}{10.9091}{mi}{65}\RG{22.634}{3.963}{7.9701}{mi}{107}\RG{26.835}{3.963}{7.9701}{rm}{43}\RG{32.819}{3.963}{7.9701}{rm}{49}\RG{40.155}{0.004}{10.9091}{rm}{61}\RG{51.489}{0.004}{10.9091}{rm}{73}\RG{55.340}{0.004}{10.9091}{rm}{109}\RG{66.119}{0.004}{10.9091}{mi}{65}\RG{74.126}{3.963}{7.9701}{mi}{107}%
}\hskip.1em\allowbreak{} 。对任意\allowbreak{}\hskip.1em \RB{24.965}{7.758}{2.722}{%
\RG{0.000}{0.004}{10.9091}{mi}{106}\RG{8.063}{0.004}{10.9091}{sy}{21}\RG{19.578}{0.004}{10.9091}{rm}{48}%
}\hskip.1em\allowbreak{} ，有

\[
\RB{234.073}{10.035}{3.236}{%
\RG{22.634}{4.506}{7.9701}{mi}{107}\RG{26.835}{4.506}{7.9701}{rm}{43}\RG{32.822}{4.506}{7.9701}{mi}{106}\RG{36.354}{4.506}{7.9701}{rm}{43}\RG{42.338}{4.506}{7.9701}{rm}{49}\RG{69.019}{4.506}{7.9701}{mi}{106}\RG{99.790}{4.506}{7.9701}{mi}{107}\RG{103.992}{4.506}{7.9701}{rm}{43}\RG{109.977}{4.506}{7.9701}{rm}{49}\RG{140.811}{4.506}{7.9701}{mi}{106}\RG{171.582}{4.506}{7.9701}{mi}{107}\RG{217.392}{4.506}{7.9701}{mi}{107}\RG{221.594}{4.506}{7.9701}{rm}{43}\RG{227.581}{4.506}{7.9701}{mi}{106}\RG{0.000}{0.002}{10.9091}{rm}{73}\RG{3.851}{0.002}{10.9091}{rm}{109}\RG{14.625}{0.002}{10.9091}{mi}{65}\RG{49.673}{0.002}{10.9091}{rm}{61}\RG{61.012}{0.002}{10.9091}{mi}{65}\RG{73.003}{0.002}{10.9091}{rm}{40}\RG{77.159}{0.002}{10.9091}{rm}{73}\RG{81.010}{0.002}{10.9091}{rm}{109}\RG{91.783}{0.002}{10.9091}{mi}{65}\RG{114.281}{0.002}{10.9091}{rm}{41}\RG{121.459}{0.002}{10.9091}{rm}{61}\RG{132.804}{0.002}{10.9091}{mi}{65}\RG{144.795}{0.002}{10.9091}{rm}{40}\RG{148.951}{0.002}{10.9091}{rm}{73}\RG{152.802}{0.002}{10.9091}{rm}{109}\RG{163.575}{0.002}{10.9091}{mi}{65}\RG{176.236}{0.002}{10.9091}{rm}{41}\RG{183.415}{0.002}{10.9091}{rm}{61}\RG{194.760}{0.002}{10.9091}{rm}{73}\RG{198.611}{0.002}{10.9091}{rm}{109}\RG{209.385}{0.002}{10.9091}{mi}{65}\RG{231.565}{0.002}{10.9091}{mi}{58}%
}
\]
因此从\allowbreak{}\hskip.1em \RB{5.965}{8.071}{0.609}{%
\RG{0.000}{0.000}{10.9091}{mi}{107}%
}\hskip.1em\allowbreak{} 开始所有相邻幂的秩均相等，即\allowbreak{}\hskip.1em \RB{64.371}{9.487}{3.238}{%
\RG{0.000}{0.000}{10.9091}{mi}{114}\RG{5.110}{0.000}{10.9091}{rm}{40}\RG{9.265}{0.000}{10.9091}{mi}{65}\RG{17.272}{3.959}{7.9701}{mi}{105}\RG{20.345}{0.000}{10.9091}{rm}{41}\RG{27.524}{0.000}{10.9091}{rm}{61}\RG{38.868}{0.000}{10.9091}{mi}{114}\RG{43.979}{0.000}{10.9091}{rm}{40}\RG{48.133}{0.000}{10.9091}{mi}{65}\RG{56.140}{3.959}{7.9701}{mi}{107}\RG{60.795}{0.000}{10.9091}{rm}{41}%
}\hskip.1em\allowbreak{} （\allowbreak{}\hskip.1em \RB{24.155}{8.071}{1.995}{%
\RG{0.000}{0.000}{10.9091}{mi}{105}\RG{6.675}{0.000}{10.9091}{sy}{21}\RG{18.190}{0.000}{10.9091}{mi}{107}%
}\hskip.1em\allowbreak{} ）。（2）序列\allowbreak{}\hskip.1em \RB{149.561}{9.275}{3.241}{%
\RG{0.000}{-0.003}{10.9091}{mi}{110}\RG{9.406}{-0.003}{10.9091}{rm}{61}\RG{20.745}{-0.003}{10.9091}{mi}{114}\RG{25.855}{-0.003}{10.9091}{rm}{40}\RG{30.010}{-0.003}{10.9091}{mi}{65}\RG{38.017}{3.956}{7.9701}{rm}{48}\RG{42.319}{-0.003}{10.9091}{rm}{41}\RG{49.504}{-0.003}{10.9091}{sy}{21}\RG{61.018}{-0.003}{10.9091}{mi}{114}\RG{66.129}{-0.003}{10.9091}{rm}{40}\RG{70.283}{-0.003}{10.9091}{mi}{65}\RG{78.290}{-0.003}{10.9091}{rm}{41}\RG{85.475}{-0.003}{10.9091}{sy}{21}\RG{96.990}{-0.003}{10.9091}{mi}{114}\RG{102.100}{-0.003}{10.9091}{rm}{40}\RG{106.254}{-0.003}{10.9091}{mi}{65}\RG{114.262}{3.956}{7.9701}{rm}{50}\RG{118.564}{-0.003}{10.9091}{rm}{41}\RG{125.748}{-0.003}{10.9091}{sy}{21}\RG{137.257}{-0.003}{10.9091}{sy}{1}\RG{142.112}{-0.003}{10.9091}{sy}{1}\RG{146.966}{-0.003}{10.9091}{sy}{1}%
}\hskip.1em\allowbreak{} 由非负整数组成。若在前\allowbreak{}\hskip.1em \RB{24.465}{7.752}{1.343}{%
\RG{0.000}{-0.003}{10.9091}{mi}{110}\RG{8.615}{-0.003}{10.9091}{rm}{43}\RG{19.165}{-0.003}{10.9091}{rm}{50}%
}\hskip.1em\allowbreak{} 项中无相邻两项相等，则\allowbreak{}\hskip.1em \RB{132.605}{9.275}{3.241}{%
\RG{0.000}{-0.003}{10.9091}{mi}{114}\RG{5.110}{-0.003}{10.9091}{rm}{40}\RG{9.264}{-0.003}{10.9091}{mi}{65}\RG{17.272}{3.956}{7.9701}{mi}{110}\RG{21.943}{3.956}{7.9701}{rm}{43}\RG{27.928}{3.956}{7.9701}{rm}{49}\RG{32.232}{-0.003}{10.9091}{rm}{41}\RG{39.587}{-0.003}{10.9091}{sy}{20}\RG{51.273}{-0.003}{10.9091}{mi}{110}\RG{60.142}{-0.003}{10.9091}{sy}{0}\RG{71.119}{-0.003}{10.9091}{rm}{40}\RG{75.273}{-0.003}{10.9091}{mi}{110}\RG{84.142}{-0.003}{10.9091}{rm}{43}\RG{94.942}{-0.003}{10.9091}{rm}{49}\RG{100.287}{-0.003}{10.9091}{rm}{41}\RG{107.640}{-0.003}{10.9091}{rm}{61}\RG{119.148}{-0.003}{10.9091}{sy}{0}\RG{127.633}{-0.003}{10.9091}{rm}{49}%
}\hskip.1em\allowbreak{} ，矛盾。故至少有一个\allowbreak{}\hskip.1em \RB{75.979}{8.679}{3.230}{%
\RG{0.000}{-0.003}{10.9091}{mi}{107}\RG{9.123}{-0.003}{10.9091}{sy}{50}\RG{19.595}{-0.003}{10.9091}{sy}{102}\RG{25.052}{-0.003}{10.9091}{rm}{48}\RG{30.393}{-0.003}{10.9091}{mi}{59}\RG{35.177}{-0.003}{10.9091}{rm}{49}\RG{40.518}{-0.003}{10.9091}{mi}{59}\RG{45.296}{-0.003}{10.9091}{mi}{58}\RG{50.085}{-0.003}{10.9091}{mi}{58}\RG{54.875}{-0.003}{10.9091}{mi}{58}\RG{59.653}{-0.003}{10.9091}{mi}{59}\RG{64.442}{-0.003}{10.9091}{mi}{110}\RG{70.821}{-0.003}{10.9091}{sy}{103}%
}\hskip.1em\allowbreak{} 使相邻秩相等，因而\allowbreak{}\hskip.1em \RB{20.769}{8.679}{3.241}{%
\RG{0.000}{-0.003}{10.9091}{mi}{115}\RG{5.032}{-0.003}{10.9091}{rm}{40}\RG{9.185}{-0.003}{10.9091}{mi}{65}\RG{17.193}{-0.003}{10.9091}{rm}{41}%
}\hskip.1em\allowbreak{} 存在且\allowbreak{}\hskip.1em \RB{62.333}{8.679}{3.241}{%
\RG{0.000}{-0.003}{10.9091}{rm}{48}\RG{8.371}{-0.003}{10.9091}{sy}{20}\RG{19.886}{-0.003}{10.9091}{mi}{115}\RG{24.917}{-0.003}{10.9091}{rm}{40}\RG{29.072}{-0.003}{10.9091}{mi}{65}\RG{37.079}{-0.003}{10.9091}{rm}{41}\RG{44.263}{-0.003}{10.9091}{sy}{20}\RG{55.778}{-0.003}{10.9091}{mi}{110}%
}\hskip.1em\allowbreak{} ；由（1），从该处起秩一直保持不变。（3）与（4）记\allowbreak{}\hskip.1em \RB{39.350}{8.285}{0.505}{%
\RG{0.000}{-0.005}{10.9091}{mi}{80}\RG{11.399}{-0.005}{10.9091}{rm}{61}\RG{22.737}{-0.005}{10.9091}{mi}{65}\RG{30.745}{-0.005}{10.9091}{mi}{66}%
}\hskip.1em\allowbreak{} 、\allowbreak{}\hskip.1em \RB{39.693}{8.285}{2.621}{%
\RG{0.000}{-0.005}{10.9091}{mi}{81}\RG{11.491}{-0.005}{10.9091}{rm}{61}\RG{22.829}{-0.005}{10.9091}{mi}{66}\RG{31.491}{-0.005}{10.9091}{mi}{65}%
}\hskip.1em\allowbreak{} ，并令\allowbreak{}\hskip.1em \RB{40.531}{8.677}{3.243}{%
\RG{0.000}{-0.005}{10.9091}{mi}{115}\RG{8.062}{-0.005}{10.9091}{rm}{61}\RG{19.400}{-0.005}{10.9091}{mi}{115}\RG{24.432}{-0.005}{10.9091}{rm}{40}\RG{28.585}{-0.005}{10.9091}{mi}{80}\RG{36.955}{-0.005}{10.9091}{rm}{41}%
}\hskip.1em\allowbreak{} 。对任意整数\allowbreak{}\hskip.1em \RB{24.415}{8.175}{1.999}{%
\RG{0.000}{-0.005}{10.9091}{mi}{96}\RG{7.513}{-0.005}{10.9091}{sy}{21}\RG{19.028}{-0.005}{10.9091}{rm}{48}%
}\hskip.1em\allowbreak{} ，有

\[
\RB{168.872}{10.107}{2.614}{%
\RG{8.369}{4.506}{7.9701}{mi}{96}\RG{11.514}{4.506}{7.9701}{rm}{43}\RG{17.498}{4.506}{7.9701}{rm}{50}\RG{52.639}{4.506}{7.9701}{mi}{96}\RG{55.785}{4.506}{7.9701}{rm}{43}\RG{61.769}{4.506}{7.9701}{rm}{49}\RG{109.803}{4.506}{7.9701}{mi}{96}\RG{112.947}{4.506}{7.9701}{rm}{43}\RG{118.932}{4.506}{7.9701}{rm}{49}\RG{154.640}{4.506}{7.9701}{mi}{96}\RG{0.000}{0.003}{10.9091}{mi}{80}\RG{24.834}{0.003}{10.9091}{rm}{61}\RG{36.172}{0.003}{10.9091}{mi}{65}\RG{44.179}{0.003}{10.9091}{mi}{81}\RG{66.074}{0.003}{10.9091}{mi}{66}\RG{74.735}{0.003}{10.9091}{mi}{59}\RG{101.343}{0.003}{10.9091}{mi}{81}\RG{126.267}{0.003}{10.9091}{rm}{61}\RG{137.605}{0.003}{10.9091}{mi}{66}\RG{146.267}{0.003}{10.9091}{mi}{80}\RG{158.237}{0.003}{10.9091}{mi}{65}\RG{166.245}{0.003}{10.9091}{mi}{59}%
}
\]
于是

\[
\RB{127.516}{10.105}{3.238}{%
\RG{17.633}{4.505}{7.9701}{mi}{96}\RG{20.777}{4.505}{7.9701}{rm}{43}\RG{26.762}{4.505}{7.9701}{rm}{50}\RG{67.491}{4.505}{7.9701}{mi}{96}\RG{70.635}{4.505}{7.9701}{rm}{43}\RG{76.620}{4.505}{7.9701}{rm}{49}\RG{117.257}{4.505}{7.9701}{mi}{96}\RG{0.000}{-0.000}{10.9091}{mi}{114}\RG{5.110}{0.001}{10.9091}{rm}{40}\RG{9.264}{-0.000}{10.9091}{mi}{80}\RG{31.067}{-0.000}{10.9091}{rm}{41}\RG{38.251}{0.001}{10.9091}{sy}{20}\RG{49.766}{-0.000}{10.9091}{mi}{114}\RG{54.876}{0.001}{10.9091}{rm}{40}\RG{59.031}{-0.000}{10.9091}{mi}{81}\RG{80.925}{-0.000}{10.9091}{rm}{41}\RG{88.109}{0.001}{10.9091}{sy}{20}\RG{99.625}{-0.000}{10.9091}{mi}{114}\RG{104.735}{0.001}{10.9091}{rm}{40}\RG{108.889}{-0.000}{10.9091}{mi}{80}\RG{120.855}{-0.000}{10.9091}{rm}{41}\RG{125.009}{0.001}{10.9091}{mi}{58}%
}
\]
当\allowbreak{}\hskip.1em \RB{24.034}{8.180}{1.995}{%
\RG{0.000}{-0.000}{10.9091}{mi}{96}\RG{7.513}{-0.000}{10.9091}{sy}{21}\RG{19.028}{-0.000}{10.9091}{mi}{115}%
}\hskip.1em\allowbreak{} 时，两端都等于\allowbreak{}\hskip.1em \RB{29.065}{8.682}{3.238}{%
\RG{0.000}{-0.000}{10.9091}{mi}{114}\RG{4.814}{-1.636}{7.9701}{sy}{49}\RG{12.965}{-0.000}{10.9091}{rm}{40}\RG{17.119}{-0.000}{10.9091}{mi}{80}\RG{25.488}{-0.000}{10.9091}{rm}{41}%
}\hskip.1em\allowbreak{} ，从而

\[
\RB{145.613}{9.799}{3.240}{%
\RG{17.724}{4.503}{7.9701}{mi}{106}\RG{0.000}{-0.002}{10.9091}{mi}{114}\RG{5.110}{-0.002}{10.9091}{rm}{40}\RG{9.264}{-0.002}{10.9091}{mi}{81}\RG{21.707}{-0.002}{10.9091}{rm}{41}\RG{28.885}{-0.002}{10.9091}{rm}{61}\RG{40.230}{-0.002}{10.9091}{mi}{114}\RG{53.195}{-0.002}{10.9091}{rm}{40}\RG{57.349}{-0.002}{10.9091}{mi}{80}\RG{65.718}{-0.002}{10.9091}{rm}{41}\RG{91.682}{-0.002}{10.9091}{rm}{40}\RG{95.845}{-0.002}{10.9091}{mi}{106}\RG{103.907}{-0.002}{10.9091}{sy}{21}\RG{115.423}{-0.002}{10.9091}{mi}{115}\RG{122.878}{-0.002}{10.9091}{rm}{43}\RG{133.613}{-0.002}{10.9091}{rm}{49}\RG{138.958}{-0.002}{10.9091}{rm}{41}\RG{143.105}{-0.002}{10.9091}{mi}{58}\RG{45.045}{-1.638}{7.9701}{sy}{49}%
}
\]
这说明\allowbreak{}\hskip.1em \RB{73.640}{8.678}{3.242}{%
\RG{0.000}{-0.004}{10.9091}{mi}{114}\RG{4.814}{-1.640}{7.9701}{sy}{49}\RG{12.965}{-0.004}{10.9091}{rm}{40}\RG{17.119}{-0.004}{10.9091}{mi}{81}\RG{25.579}{-0.004}{10.9091}{rm}{41}\RG{32.997}{-0.004}{10.9091}{rm}{61}\RG{44.576}{-0.004}{10.9091}{mi}{114}\RG{49.390}{-1.640}{7.9701}{sy}{49}\RG{57.541}{-0.004}{10.9091}{rm}{40}\RG{61.695}{-0.004}{10.9091}{mi}{80}\RG{70.064}{-0.004}{10.9091}{rm}{41}%
}\hskip.1em\allowbreak{} ，同时\allowbreak{}\hskip.1em \RB{76.845}{8.678}{3.242}{%
\RG{0.000}{-0.004}{10.9091}{mi}{115}\RG{5.032}{-0.004}{10.9091}{rm}{40}\RG{9.185}{-0.004}{10.9091}{mi}{81}\RG{17.645}{-0.004}{10.9091}{rm}{41}\RG{25.067}{-0.004}{10.9091}{sy}{20}\RG{36.820}{-0.004}{10.9091}{mi}{115}\RG{41.852}{-0.004}{10.9091}{rm}{40}\RG{46.005}{-0.004}{10.9091}{mi}{80}\RG{54.374}{-0.004}{10.9091}{rm}{41}\RG{61.039}{-0.004}{10.9091}{rm}{43}\RG{71.872}{-0.004}{10.9091}{rm}{49}%
}\hskip.1em\allowbreak{} 。交换\allowbreak{}\hskip.1em \RB{20.376}{8.176}{2.620}{%
\RG{0.000}{-0.004}{10.9091}{mi}{80}\RG{7.178}{-0.004}{10.9091}{mi}{59}\RG{11.956}{-0.004}{10.9091}{mi}{81}%
}\hskip.1em\allowbreak{} ，又得\allowbreak{}\hskip.1em \RB{76.845}{8.678}{3.242}{%
\RG{0.000}{-0.004}{10.9091}{mi}{115}\RG{5.032}{-0.004}{10.9091}{rm}{40}\RG{9.185}{-0.004}{10.9091}{mi}{80}\RG{17.555}{-0.004}{10.9091}{rm}{41}\RG{24.975}{-0.004}{10.9091}{sy}{20}\RG{36.728}{-0.004}{10.9091}{mi}{115}\RG{41.760}{-0.004}{10.9091}{rm}{40}\RG{45.914}{-0.004}{10.9091}{mi}{81}\RG{54.374}{-0.004}{10.9091}{rm}{41}\RG{61.039}{-0.004}{10.9091}{rm}{43}\RG{71.872}{-0.004}{10.9091}{rm}{49}%
}\hskip.1em\allowbreak{} ，从而\allowbreak{}\hskip.1em \RB{82.419}{8.678}{3.242}{%
\RG{0.000}{-0.004}{10.9091}{sy}{106}\RG{3.030}{-0.004}{10.9091}{mi}{115}\RG{8.062}{-0.004}{10.9091}{rm}{40}\RG{12.215}{-0.004}{10.9091}{mi}{80}\RG{20.585}{-0.004}{10.9091}{rm}{41}\RG{27.163}{-0.004}{10.9091}{sy}{0}\RG{38.072}{-0.004}{10.9091}{mi}{115}\RG{43.103}{-0.004}{10.9091}{rm}{40}\RG{47.257}{-0.004}{10.9091}{mi}{81}\RG{55.717}{-0.004}{10.9091}{rm}{41}\RG{59.871}{-0.004}{10.9091}{sy}{106}\RG{65.925}{-0.004}{10.9091}{sy}{20}\RG{77.446}{-0.004}{10.9091}{rm}{49}%
}\hskip.1em\allowbreak{} ，两问均得证。取二阶矩阵

\[
\RB{149.334}{19.520}{14.068}{%
\RG{22.376}{18.665}{9.0569}{ex}{48}\RG{50.040}{18.665}{9.0569}{ex}{49}\RG{109.420}{18.666}{9.0569}{ex}{48}\RG{137.084}{18.666}{9.0569}{ex}{49}\RG{30.301}{7.806}{10.9091}{rm}{49}\RG{44.701}{7.806}{10.9091}{rm}{48}\RG{117.345}{7.806}{10.9091}{rm}{48}\RG{131.745}{7.806}{10.9091}{rm}{49}\RG{22.376}{2.364}{9.0569}{ex}{64}\RG{50.040}{2.364}{9.0569}{ex}{65}\RG{109.420}{2.364}{9.0569}{ex}{64}\RG{137.084}{2.364}{9.0569}{ex}{65}\RG{0.000}{-0.002}{10.9091}{mi}{65}\RG{11.037}{-0.002}{10.9091}{rm}{61}\RG{59.783}{-0.002}{10.9091}{mi}{59}\RG{86.379}{-0.002}{10.9091}{mi}{66}\RG{98.082}{-0.002}{10.9091}{rm}{61}\RG{146.826}{-0.002}{10.9091}{mi}{58}\RG{30.301}{-9.130}{10.9091}{rm}{48}\RG{44.701}{-9.130}{10.9091}{rm}{48}\RG{117.345}{-9.130}{10.9091}{rm}{48}\RG{131.745}{-9.130}{10.9091}{rm}{48}%
}
\]
则\allowbreak{}\hskip.1em \RB{39.647}{8.288}{0.501}{%
\RG{0.000}{-0.001}{10.9091}{mi}{65}\RG{8.007}{-0.001}{10.9091}{mi}{66}\RG{19.703}{-0.001}{10.9091}{rm}{61}\RG{31.042}{-0.001}{10.9091}{mi}{66}%
}\hskip.1em\allowbreak{} 、\allowbreak{}\hskip.1em \RB{39.461}{8.288}{0.730}{%
\RG{0.000}{-0.001}{10.9091}{mi}{66}\RG{8.662}{-0.001}{10.9091}{mi}{65}\RG{19.703}{-0.001}{10.9091}{rm}{61}\RG{31.041}{-0.001}{10.9091}{mi}{79}%
}\hskip.1em\allowbreak{} 。由于\allowbreak{}\hskip.1em \RB{31.455}{8.310}{2.857}{%
\RG{0.000}{-0.001}{10.9091}{mi}{66}\RG{11.695}{-0.001}{10.9091}{sy}{54}\RG{11.695}{-0.001}{10.9091}{rm}{61}\RG{23.035}{-0.001}{10.9091}{mi}{79}%
}\hskip.1em\allowbreak{} 且\allowbreak{}\hskip.1em \RB{35.755}{9.276}{0.730}{%
\RG{0.000}{-0.001}{10.9091}{mi}{66}\RG{8.665}{3.958}{7.9701}{rm}{50}\RG{15.997}{-0.001}{10.9091}{rm}{61}\RG{27.335}{-0.001}{10.9091}{mi}{79}%
}\hskip.1em\allowbreak{} ，两者从零次幂开始的秩序列分别为

\[
\RB{123.835}{25.051}{23.196}{%
\RG{19.930}{17.296}{10.9091}{mi}{106}\RG{53.950}{17.296}{10.9091}{rm}{48}\RG{68.350}{17.296}{10.9091}{rm}{49}\RG{82.750}{17.296}{10.9091}{rm}{50}\RG{97.150}{17.296}{10.9091}{rm}{51}\RG{111.542}{17.296}{10.9091}{sy}{1}\RG{116.385}{17.296}{10.9091}{sy}{1}\RG{121.240}{17.296}{10.9091}{sy}{1}\RG{5.111}{10.515}{9.0569}{ex}{16}\RG{39.484}{10.515}{9.0569}{ex}{17}\RG{35.500}{3.957}{7.9701}{mi}{106}\RG{0.000}{-0.002}{10.9091}{mi}{114}\RG{10.520}{-0.002}{10.9091}{rm}{40}\RG{14.674}{-0.002}{10.9091}{mi}{65}\RG{22.681}{-0.002}{10.9091}{mi}{66}\RG{31.346}{-0.002}{10.9091}{rm}{41}\RG{53.950}{-0.002}{10.9091}{rm}{50}\RG{68.350}{-0.002}{10.9091}{rm}{49}\RG{82.750}{-0.002}{10.9091}{rm}{48}\RG{97.150}{-0.002}{10.9091}{rm}{48}\RG{111.542}{-0.002}{10.9091}{sy}{1}\RG{116.385}{-0.002}{10.9091}{sy}{1}\RG{121.240}{-0.002}{10.9091}{sy}{1}\RG{5.111}{-6.765}{9.0569}{ex}{16}\RG{39.484}{-6.765}{9.0569}{ex}{17}\RG{35.500}{-13.323}{7.9701}{mi}{106}\RG{0.000}{-17.281}{10.9091}{mi}{114}\RG{10.520}{-17.281}{10.9091}{rm}{40}\RG{14.674}{-17.281}{10.9091}{mi}{66}\RG{23.335}{-17.281}{10.9091}{mi}{65}\RG{31.346}{-17.281}{10.9091}{rm}{41}\RG{53.950}{-17.281}{10.9091}{rm}{50}\RG{68.350}{-17.281}{10.9091}{rm}{48}\RG{82.750}{-17.281}{10.9091}{rm}{48}\RG{97.150}{-17.281}{10.9091}{rm}{48}\RG{111.542}{-17.281}{10.9091}{sy}{1}\RG{116.385}{-17.281}{10.9091}{sy}{1}\RG{121.240}{-17.281}{10.9091}{sy}{1}\RR{-4.528}{12.035}{133.325}{0.362}%
}
\]
故\allowbreak{}\hskip.1em \RB{50.784}{8.680}{3.240}{%
\RG{0.000}{-0.002}{10.9091}{mi}{115}\RG{5.031}{-0.002}{10.9091}{rm}{40}\RG{9.185}{-0.002}{10.9091}{mi}{65}\RG{17.193}{-0.002}{10.9091}{mi}{66}\RG{25.858}{-0.002}{10.9091}{rm}{41}\RG{33.593}{-0.002}{10.9091}{rm}{61}\RG{45.484}{-0.002}{10.9091}{rm}{50}%
}\hskip.1em\allowbreak{} 、\allowbreak{}\hskip.1em \RB{50.456}{8.680}{3.240}{%
\RG{0.000}{-0.002}{10.9091}{mi}{115}\RG{5.032}{-0.002}{10.9091}{rm}{40}\RG{9.185}{-0.002}{10.9091}{mi}{66}\RG{17.847}{-0.002}{10.9091}{mi}{65}\RG{25.858}{-0.002}{10.9091}{rm}{41}\RG{33.593}{-0.002}{10.9091}{rm}{61}\RG{45.484}{-0.002}{10.9091}{rm}{49}%
}\hskip.1em\allowbreak{} ，其差的绝对值等于\allowbreak{}\hskip.1em \RB{4.973}{7.753}{0.502}{%
\RG{0.000}{-0.002}{10.9091}{rm}{49}%
}\hskip.1em\allowbreak{} 。任意\allowbreak{}\hskip.1em \RB{27.143}{7.753}{0.807}{%
\RG{0.000}{-0.002}{10.9091}{mi}{110}\RG{9.960}{-0.002}{10.9091}{mi}{62}\RG{21.843}{-0.002}{10.9091}{rm}{50}%
}\hskip.1em\allowbreak{} 可在二者右下角同添\allowbreak{}\hskip.1em \RB{19.326}{7.949}{2.211}{%
\RG{0.000}{-0.002}{10.9091}{mi}{73}\RG{4.711}{-1.638}{7.9701}{mi}{110}\RG{9.383}{-1.638}{7.9701}{sy}{0}\RG{15.370}{-1.638}{7.9701}{rm}{50}%
}\hskip.1em\allowbreak{} ；\allowbreak{}\hskip.1em \RB{25.714}{7.753}{0.611}{%
\RG{0.000}{-0.002}{10.9091}{mi}{110}\RG{9.406}{-0.002}{10.9091}{rm}{61}\RG{20.741}{-0.002}{10.9091}{rm}{49}%
}\hskip.1em\allowbreak{} 时因\allowbreak{}\hskip.1em \RB{47.905}{8.287}{0.502}{%
\RG{0.000}{-0.002}{10.9091}{mi}{65}\RG{8.007}{-0.002}{10.9091}{mi}{66}\RG{19.703}{-0.002}{10.9091}{rm}{61}\RG{31.041}{-0.002}{10.9091}{mi}{66}\RG{39.703}{-0.002}{10.9091}{mi}{65}%
}\hskip.1em\allowbreak{} ，稳定指数之差恒为\allowbreak{}\hskip.1em \RB{5.387}{7.753}{0.731}{%
\RG{0.000}{-0.002}{10.9091}{rm}{48}%
}\hskip.1em\allowbreak{} 。\allowbreak{}\hskip.1em \RB{8.365}{7.993}{0.502}{%
\RG{0.000}{-0.002}{10.9091}{am}{3}%
}\hskip.1em\allowbreak{}

小结：幂的像空间逐层包含，一旦相邻维数相等便永久稳定。比较\allowbreak{}\hskip.1em \RB{16.613}{8.285}{0.505}{%
\RG{0.000}{-0.005}{10.9091}{mi}{65}\RG{8.007}{-0.005}{10.9091}{mi}{66}%
}\hskip.1em\allowbreak{} 与\allowbreak{}\hskip.1em \RB{16.864}{8.285}{0.505}{%
\RG{0.000}{-0.005}{10.9091}{mi}{66}\RG{8.662}{-0.005}{10.9091}{mi}{65}%
}\hskip.1em\allowbreak{} 时，用相差一次幂的乘积恒等式夹住秩，即可同时控制稳定秩与稳定指数。原书勘误.原书第22页上述例子写作\allowbreak{}\hskip.1em \RB{44.001}{7.661}{2.885}{%
\RG{0.000}{0.002}{7.9701}{mi}{115}\RG{4.474}{0.002}{7.9701}{rm}{40}\RG{8.186}{0.002}{7.9701}{mi}{65}\RG{15.345}{0.002}{7.9701}{mi}{66}\RG{23.065}{0.002}{7.9701}{rm}{41}\RG{29.422}{0.002}{7.9701}{rm}{61}\RG{39.502}{0.002}{7.9701}{rm}{49}%
}\hskip.1em\allowbreak{} 、\allowbreak{}\hskip.1em \RB{44.393}{7.661}{2.885}{%
\RG{0.000}{0.002}{7.9701}{mi}{115}\RG{4.474}{0.002}{7.9701}{rm}{40}\RG{8.186}{0.002}{7.9701}{mi}{66}\RG{15.909}{0.002}{7.9701}{mi}{65}\RG{23.065}{0.002}{7.9701}{rm}{41}\RG{29.422}{0.002}{7.9701}{rm}{61}\RG{39.502}{0.002}{7.9701}{rm}{48}%
}\hskip.1em\allowbreak{} ；按本题\allowbreak{}\hskip.1em \RB{32.817}{8.176}{1.993}{%
\RG{0.000}{0.002}{7.9701}{mi}{65}\RG{7.159}{3.466}{7.9701}{rm}{48}\RG{13.873}{0.002}{7.9701}{rm}{61}\RG{23.949}{0.002}{7.9701}{mi}{73}\RG{28.145}{-1.430}{7.9701}{mi}{110}%
}\hskip.1em\allowbreak{} 及稳定指数的定义，正确数值应为\allowbreak{}\hskip.1em \RB{4.786}{6.859}{0.500}{%
\RG{0.000}{0.002}{7.9701}{rm}{50}%
}\hskip.1em\allowbreak{} 、\allowbreak{}\hskip.1em \RB{4.500}{6.859}{0.500}{%
\RG{0.000}{0.002}{7.9701}{rm}{49}%
}\hskip.1em\allowbreak{} 。

% Original page 4.
\section*{2015A06}
设\allowbreak{}\hskip.1em \RB{42.319}{8.680}{3.616}{%
\RG{0.000}{-0.002}{10.9091}{mi}{65}\RG{11.038}{-0.002}{10.9091}{rm}{61}\RG{22.373}{-0.002}{10.9091}{rm}{40}\RG{26.530}{-0.002}{10.9091}{mi}{97}\RG{32.137}{-1.638}{7.9701}{mi}{105}\RG{34.760}{-1.638}{7.9701}{mi}{106}\RG{38.743}{-0.002}{10.9091}{rm}{41}%
}\hskip.1em\allowbreak{} 是\allowbreak{}\hskip.1em \RB{6.555}{5.309}{0.611}{%
\RG{0.000}{-0.002}{10.9091}{mi}{110}%
}\hskip.1em\allowbreak{} 阶方阵，\allowbreak{}\hskip.1em \RB{14.123}{8.287}{3.616}{%
\RG{0.000}{-0.002}{10.9091}{mi}{65}\RG{8.007}{-1.638}{7.9701}{mi}{105}\RG{10.630}{-1.638}{7.9701}{mi}{106}%
}\hskip.1em\allowbreak{} 表示元素\allowbreak{}\hskip.1em \RB{11.723}{5.309}{3.616}{%
\RG{0.000}{-0.002}{10.9091}{mi}{97}\RG{5.607}{-1.638}{7.9701}{mi}{105}\RG{8.230}{-1.638}{7.9701}{mi}{106}%
}\hskip.1em\allowbreak{} 对应的代数余子式。设

\[
\RB{224.338}{7.752}{2.728}{%
\RG{0.000}{-0.003}{10.9091}{rm}{49}\RG{8.371}{-0.003}{10.9091}{sy}{20}\RG{19.886}{-0.003}{10.9091}{mi}{105}\RG{30.863}{-0.003}{10.9091}{mi}{60}\RG{42.201}{-0.003}{10.9091}{sy}{1}\RG{47.045}{-0.003}{10.9091}{sy}{1}\RG{51.899}{-0.003}{10.9091}{sy}{1}\RG{57.958}{-0.003}{10.9091}{mi}{60}\RG{69.293}{-0.003}{10.9091}{mi}{105}\RG{80.112}{-0.003}{10.9091}{sy}{20}\RG{91.626}{-0.003}{10.9091}{mi}{110}\RG{98.008}{-0.003}{10.9091}{mi}{59}\RG{124.606}{-0.003}{10.9091}{rm}{49}\RG{132.977}{-0.003}{10.9091}{sy}{20}\RG{144.493}{-0.003}{10.9091}{mi}{106}\RG{156.245}{-0.003}{10.9091}{mi}{60}\RG{167.583}{-0.003}{10.9091}{sy}{1}\RG{172.426}{-0.003}{10.9091}{sy}{1}\RG{177.281}{-0.003}{10.9091}{sy}{1}\RG{183.340}{-0.003}{10.9091}{mi}{60}\RG{194.675}{-0.003}{10.9091}{mi}{106}\RG{206.268}{-0.003}{10.9091}{sy}{20}\RG{217.784}{-0.003}{10.9091}{mi}{110}\RG{23.530}{-1.639}{7.9701}{rm}{49}\RG{72.941}{-1.639}{7.9701}{mi}{114}\RG{148.912}{-1.639}{7.9701}{rm}{49}\RG{199.097}{-1.639}{7.9701}{mi}{114}%
}
\]
为两组给定的指标，证明：

\[
\RB{309.464}{28.211}{38.432}{%
\RG{83.285}{27.521}{9.0569}{ex}{12}\RG{193.272}{27.521}{9.0569}{ex}{12}\RG{83.285}{22.087}{9.0569}{ex}{12}\RG{193.272}{22.087}{9.0569}{ex}{12}\RG{86.307}{17.875}{10.9091}{mi}{65}\RG{117.113}{17.875}{10.9091}{mi}{65}\RG{147.915}{17.875}{10.9091}{sy}{1}\RG{152.758}{17.875}{10.9091}{sy}{1}\RG{157.613}{17.875}{10.9091}{sy}{1}\RG{171.520}{17.875}{10.9091}{mi}{65}\RG{83.285}{16.653}{9.0569}{ex}{12}\RG{193.272}{16.653}{9.0569}{ex}{12}\RG{94.315}{16.238}{7.9701}{mi}{105}\RG{100.710}{16.238}{7.9701}{mi}{106}\RG{125.120}{16.238}{7.9701}{mi}{105}\RG{131.515}{16.238}{7.9701}{mi}{106}\RG{179.528}{16.238}{7.9701}{mi}{105}\RG{185.929}{16.238}{7.9701}{mi}{106}\RG{182.149}{15.332}{5.9776}{mi}{114}\RG{96.935}{15.232}{5.9776}{rm}{49}\RG{103.823}{15.232}{5.9776}{rm}{49}\RG{127.741}{15.232}{5.9776}{rm}{50}\RG{134.628}{15.232}{5.9776}{rm}{49}\RG{189.043}{15.232}{5.9776}{rm}{49}\RG{83.285}{11.218}{9.0569}{ex}{12}\RG{193.272}{11.218}{9.0569}{ex}{12}\RG{12.309}{10.856}{9.0569}{ex}{48}\RG{60.992}{10.856}{9.0569}{ex}{49}\RG{218.667}{10.856}{9.0569}{ex}{48}\RG{267.349}{10.856}{9.0569}{ex}{49}\RG{83.285}{5.785}{9.0569}{ex}{12}\RG{193.272}{5.785}{9.0569}{ex}{12}\RG{86.307}{0.938}{10.9091}{mi}{65}\RG{117.113}{0.938}{10.9091}{mi}{65}\RG{147.915}{0.938}{10.9091}{sy}{1}\RG{152.758}{0.938}{10.9091}{sy}{1}\RG{157.613}{0.938}{10.9091}{sy}{1}\RG{171.520}{0.938}{10.9091}{mi}{65}\RG{83.285}{0.350}{9.0569}{ex}{12}\RG{193.272}{0.350}{9.0569}{ex}{12}\RG{20.235}{-0.004}{10.9091}{mi}{106}\RG{34.330}{-0.004}{10.9091}{sy}{1}\RG{39.174}{-0.004}{10.9091}{sy}{1}\RG{44.028}{-0.004}{10.9091}{sy}{1}\RG{52.432}{-0.004}{10.9091}{mi}{106}\RG{227.367}{-0.004}{10.9091}{mi}{105}\RG{240.687}{-0.004}{10.9091}{sy}{1}\RG{245.531}{-0.004}{10.9091}{sy}{1}\RG{250.385}{-0.004}{10.9091}{sy}{1}\RG{258.789}{-0.004}{10.9091}{mi}{105}\RG{94.315}{-0.698}{7.9701}{mi}{105}\RG{100.709}{-0.698}{7.9701}{mi}{106}\RG{125.120}{-0.698}{7.9701}{mi}{105}\RG{131.515}{-0.698}{7.9701}{mi}{106}\RG{179.528}{-0.698}{7.9701}{mi}{105}\RG{185.929}{-0.698}{7.9701}{mi}{106}\RG{24.654}{-1.640}{7.9701}{rm}{49}\RG{56.851}{-1.640}{7.9701}{mi}{114}\RG{182.149}{-1.604}{5.9776}{mi}{114}\RG{231.011}{-1.640}{7.9701}{rm}{49}\RG{262.434}{-1.640}{7.9701}{mi}{114}\RG{96.935}{-1.705}{5.9776}{rm}{49}\RG{103.823}{-1.705}{5.9776}{rm}{50}\RG{127.741}{-1.705}{5.9776}{rm}{50}\RG{134.628}{-1.705}{5.9776}{rm}{50}\RG{189.042}{-1.705}{5.9776}{rm}{50}\RG{8.007}{-3.307}{7.9701}{sy}{3}\RG{292.978}{-3.307}{7.9701}{mi}{114}\RG{296.666}{-3.307}{7.9701}{sy}{0}\RG{302.654}{-3.307}{7.9701}{rm}{49}\RG{213.631}{-4.256}{9.0569}{ex}{98}\RG{83.285}{-5.085}{9.0569}{ex}{12}\RG{193.272}{-5.085}{9.0569}{ex}{12}\RG{12.309}{-5.446}{9.0569}{ex}{64}\RG{60.992}{-5.446}{9.0569}{ex}{65}\RG{218.667}{-5.446}{9.0569}{ex}{64}\RG{267.349}{-5.446}{9.0569}{ex}{65}\RG{0.000}{-7.812}{10.9091}{mi}{65}\RG{71.946}{-7.812}{10.9091}{rm}{61}\RG{199.321}{-7.812}{10.9091}{rm}{61}\RG{210.660}{-7.812}{10.9091}{mi}{65}\RG{278.910}{-7.812}{10.9091}{sy}{106}\RG{281.941}{-7.812}{10.9091}{mi}{65}\RG{289.948}{-7.812}{10.9091}{sy}{106}\RG{306.956}{-7.812}{10.9091}{mi}{58}\RG{83.285}{-10.518}{9.0569}{ex}{12}\RG{193.272}{-10.518}{9.0569}{ex}{12}\RT{95.695}{-10.635}{10.9091}{.}\RT{126.500}{-10.635}{10.9091}{.}\RT{180.911}{-10.635}{10.9091}{.}\RT{149.525}{-11.541}{10.9091}{.}\RT{95.695}{-14.258}{10.9091}{.}\RT{126.500}{-14.258}{10.9091}{.}\RT{153.704}{-14.258}{10.9091}{.}\RT{180.911}{-14.258}{10.9091}{.}\RG{83.285}{-15.953}{9.0569}{ex}{12}\RG{193.272}{-15.953}{9.0569}{ex}{12}\RG{21.009}{-16.939}{10.9091}{mi}{105}\RG{34.330}{-16.939}{10.9091}{sy}{1}\RG{39.174}{-16.939}{10.9091}{sy}{1}\RG{44.028}{-16.939}{10.9091}{sy}{1}\RG{52.432}{-16.939}{10.9091}{mi}{105}\RG{226.592}{-16.940}{10.9091}{mi}{106}\RG{240.687}{-16.940}{10.9091}{sy}{1}\RG{245.531}{-16.940}{10.9091}{sy}{1}\RG{250.385}{-16.940}{10.9091}{sy}{1}\RG{258.789}{-16.940}{10.9091}{mi}{106}\RT{157.884}{-16.975}{10.9091}{.}\RT{95.695}{-17.881}{10.9091}{.}\RT{126.500}{-17.881}{10.9091}{.}\RT{180.911}{-17.881}{10.9091}{.}\RG{24.654}{-18.575}{7.9701}{rm}{49}\RG{56.075}{-18.575}{7.9701}{mi}{114}\RG{231.011}{-18.576}{7.9701}{rm}{49}\RG{263.208}{-18.576}{7.9701}{mi}{114}\RG{83.285}{-21.387}{9.0569}{ex}{12}\RG{193.272}{-21.387}{9.0569}{ex}{12}\RG{83.285}{-26.821}{9.0569}{ex}{12}\RG{193.272}{-26.821}{9.0569}{ex}{12}\RG{83.285}{-32.255}{9.0569}{ex}{12}\RG{193.272}{-32.255}{9.0569}{ex}{12}\RG{86.305}{-34.817}{10.9091}{mi}{65}\RG{117.109}{-34.817}{10.9091}{mi}{65}\RG{147.915}{-34.817}{10.9091}{sy}{1}\RG{152.758}{-34.817}{10.9091}{sy}{1}\RG{157.613}{-34.817}{10.9091}{sy}{1}\RG{171.517}{-34.817}{10.9091}{mi}{65}\RG{94.312}{-36.454}{7.9701}{mi}{105}\RG{100.706}{-36.454}{7.9701}{mi}{106}\RG{125.117}{-36.454}{7.9701}{mi}{105}\RG{131.512}{-36.454}{7.9701}{mi}{106}\RG{179.525}{-36.454}{7.9701}{mi}{105}\RG{185.925}{-36.454}{7.9701}{mi}{106}\RG{103.820}{-37.359}{5.9776}{mi}{114}\RG{134.625}{-37.359}{5.9776}{mi}{114}\RG{182.145}{-37.359}{5.9776}{mi}{114}\RG{189.039}{-37.359}{5.9776}{mi}{114}\RG{96.933}{-37.460}{5.9776}{rm}{49}\RG{127.738}{-37.460}{5.9776}{rm}{50}%
}
\]
记\allowbreak{}\hskip.1em \RB{70.207}{8.678}{3.231}{%
\RG{0.000}{-0.004}{10.9091}{mi}{73}\RG{8.599}{-0.004}{10.9091}{rm}{61}\RG{19.937}{-0.004}{10.9091}{sy}{102}\RG{25.392}{-0.004}{10.9091}{mi}{105}\RG{29.035}{-1.640}{7.9701}{rm}{49}\RG{33.338}{-0.004}{10.9091}{mi}{59}\RG{38.116}{-0.004}{10.9091}{mi}{58}\RG{42.905}{-0.004}{10.9091}{mi}{58}\RG{47.684}{-0.004}{10.9091}{mi}{58}\RG{52.473}{-0.004}{10.9091}{mi}{59}\RG{57.262}{-0.004}{10.9091}{mi}{105}\RG{60.908}{-1.640}{7.9701}{mi}{114}\RG{65.049}{-0.004}{10.9091}{sy}{103}%
}\hskip.1em\allowbreak{} 、\allowbreak{}\hskip.1em \RB{73.103}{8.678}{3.231}{%
\RG{0.000}{-0.004}{10.9091}{mi}{74}\RG{9.944}{-0.004}{10.9091}{rm}{61}\RG{21.283}{-0.004}{10.9091}{sy}{102}\RG{26.737}{-0.004}{10.9091}{mi}{106}\RG{31.156}{-1.640}{7.9701}{rm}{49}\RG{35.458}{-0.004}{10.9091}{mi}{59}\RG{40.236}{-0.004}{10.9091}{mi}{58}\RG{45.025}{-0.004}{10.9091}{mi}{58}\RG{49.804}{-0.004}{10.9091}{mi}{58}\RG{54.593}{-0.004}{10.9091}{mi}{59}\RG{59.382}{-0.004}{10.9091}{mi}{106}\RG{63.804}{-1.640}{7.9701}{mi}{114}\RG{67.945}{-0.004}{10.9091}{sy}{103}%
}\hskip.1em\allowbreak{} ，\allowbreak{}\hskip.1em \RB{24.896}{7.947}{2.609}{%
\RG{0.000}{-0.004}{10.9091}{mi}{73}\RG{5.569}{3.955}{7.9701}{mi}{99}\RG{9.356}{-0.004}{10.9091}{mi}{59}\RG{14.135}{-0.004}{10.9091}{mi}{74}\RG{21.055}{3.955}{7.9701}{mi}{99}%
}\hskip.1em\allowbreak{} 为相应补集，\allowbreak{}\hskip.1em \RB{61.783}{8.678}{3.242}{%
\RG{0.000}{-0.004}{10.9091}{mi}{27}\RG{6.463}{-0.004}{10.9091}{rm}{40}\RG{10.616}{-0.004}{10.9091}{mi}{73}\RG{16.185}{-0.004}{10.9091}{rm}{41}\RG{23.363}{-0.004}{10.9091}{rm}{61}\RG{34.707}{7.253}{9.0569}{ex}{80}\RG{44.268}{-2.257}{7.9701}{mi}{105}\RG{46.889}{-2.257}{7.9701}{sy}{50}\RG{52.021}{-2.257}{7.9701}{mi}{73}\RG{58.195}{-0.004}{10.9091}{mi}{105}%
}\hskip.1em\allowbreak{} 。原题左端表示\allowbreak{}\hskip.1em \RB{11.847}{8.285}{0.504}{%
\RG{0.000}{-0.004}{10.9091}{mi}{65}\RG{8.007}{3.955}{7.9701}{sy}{3}%
}\hskip.1em\allowbreak{}的相应子式；右端的\allowbreak{}\hskip.1em \RB{8.553}{10.790}{0.504}{%
\RG{0.000}{-0.004}{10.9091}{mi}{65}\RG{2.972}{3.552}{9.0569}{ex}{98}%
}\hskip.1em\allowbreak{} 表示代数余子式，其值为\allowbreak{}\hskip.1em \RB{115.201}{9.890}{3.242}{%
\RG{0.000}{-0.004}{10.9091}{rm}{40}\RG{4.154}{-0.004}{10.9091}{sy}{0}\RG{12.639}{-0.004}{10.9091}{rm}{49}\RG{17.985}{-0.004}{10.9091}{rm}{41}\RG{22.134}{3.955}{7.9701}{mi}{27}\RG{26.793}{3.955}{7.9701}{rm}{40}\RG{29.787}{3.955}{7.9701}{mi}{73}\RG{33.690}{3.955}{7.9701}{rm}{41}\RG{36.682}{3.955}{7.9701}{rm}{43}\RG{42.672}{3.955}{7.9701}{mi}{27}\RG{47.331}{3.955}{7.9701}{rm}{40}\RG{50.325}{3.955}{7.9701}{mi}{74}\RG{55.211}{3.955}{7.9701}{rm}{41}\RG{60.475}{-0.004}{10.9091}{rm}{100}\RG{66.410}{-0.004}{10.9091}{rm}{101}\RG{71.155}{-0.004}{10.9091}{rm}{116}\RG{77.130}{-0.004}{10.9091}{mi}{65}\RG{85.137}{-0.004}{10.9091}{rm}{91}\RG{88.167}{-0.004}{10.9091}{mi}{73}\RG{93.736}{3.955}{7.9701}{mi}{99}\RG{97.524}{-0.004}{10.9091}{mi}{59}\RG{102.302}{-0.004}{10.9091}{mi}{74}\RG{109.223}{3.955}{7.9701}{mi}{99}\RG{113.010}{-0.004}{10.9091}{rm}{93}%
}\hskip.1em\allowbreak{} 。当\allowbreak{}\hskip.1em \RB{27.729}{5.307}{0.613}{%
\RG{0.000}{-0.004}{10.9091}{mi}{114}\RG{8.989}{-0.004}{10.9091}{rm}{61}\RG{21.175}{-0.004}{10.9091}{mi}{110}%
}\hskip.1em\allowbreak{} 时，约定零阶行列式为\allowbreak{}\hskip.1em \RB{4.973}{7.751}{0.504}{%
\RG{0.000}{-0.004}{10.9091}{rm}{49}%
}\hskip.1em\allowbreak{} 。由于\allowbreak{}\hskip.1em \RB{55.714}{8.678}{3.618}{%
\RG{0.000}{-0.004}{10.9091}{rm}{40}\RG{4.155}{-0.004}{10.9091}{mi}{65}\RG{12.162}{3.955}{7.9701}{sy}{3}\RG{16.464}{-0.004}{10.9091}{rm}{41}\RG{20.617}{-1.640}{7.9701}{mi}{106}\RG{24.146}{-1.640}{7.9701}{mi}{105}\RG{30.253}{-0.004}{10.9091}{rm}{61}\RG{41.591}{-0.004}{10.9091}{mi}{65}\RG{49.599}{-1.640}{7.9701}{mi}{105}\RG{52.222}{-1.640}{7.9701}{mi}{106}%
}\hskip.1em\allowbreak{} ，原题的第一个等号由定义即得；以下证明第二个等号。

证明：\allowbreak{}\hskip.1em \RB{24.448}{7.756}{0.607}{%
\RG{0.000}{0.002}{10.9091}{mi}{114}\RG{8.141}{0.002}{10.9091}{rm}{61}\RG{19.475}{0.002}{10.9091}{rm}{49}%
}\hskip.1em\allowbreak{} 时，结论就是代数余子式的定义。以下设\allowbreak{}\hskip.1em \RB{24.955}{7.756}{1.993}{%
\RG{0.000}{0.002}{10.9091}{mi}{114}\RG{8.141}{0.002}{10.9091}{sy}{21}\RG{19.655}{0.002}{10.9091}{rm}{50}%
}\hskip.1em\allowbreak{} ，记\allowbreak{}\hskip.1em \RB{48.390}{8.291}{0.607}{%
\RG{0.000}{0.002}{10.9091}{mi}{68}\RG{12.196}{0.002}{10.9091}{rm}{61}\RG{23.531}{0.002}{10.9091}{rm}{100}\RG{29.465}{0.002}{10.9091}{rm}{101}\RG{34.211}{0.002}{10.9091}{rm}{116}\RG{40.188}{0.002}{10.9091}{mi}{65}%
}\hskip.1em\allowbreak{} 。先设\allowbreak{}\hskip.1em \RB{29.071}{8.308}{2.859}{%
\RG{0.000}{-0.003}{10.9091}{mi}{68}\RG{12.273}{-0.003}{10.9091}{sy}{54}\RG{12.273}{-0.003}{10.9091}{rm}{61}\RG{23.684}{-0.003}{10.9091}{rm}{48}%
}\hskip.1em\allowbreak{} 。将\allowbreak{}\hskip.1em \RB{8.202}{8.286}{0.503}{%
\RG{0.000}{-0.003}{10.9091}{mi}{65}%
}\hskip.1em\allowbreak{} 的第\allowbreak{}\hskip.1em \RB{7.911}{7.665}{3.545}{%
\RG{0.000}{-0.003}{10.9091}{mi}{105}\RG{3.644}{-1.639}{7.9701}{mi}{112}%
}\hskip.1em\allowbreak{} 行换为第\allowbreak{}\hskip.1em \RB{8.687}{7.665}{3.545}{%
\RG{0.000}{-0.003}{10.9091}{mi}{106}\RG{4.419}{-1.639}{7.9701}{mi}{112}%
}\hskip.1em\allowbreak{} 个标准单位行向量\allowbreak{}\hskip.1em \RB{11.820}{8.507}{5.159}{%
\RG{0.000}{-0.003}{10.9091}{mi}{101}\RG{4.951}{3.956}{7.9701}{sy}{48}\RG{4.951}{-2.699}{7.9701}{mi}{106}\RG{8.065}{-3.605}{5.9776}{mi}{112}%
}\hskip.1em\allowbreak{} （\allowbreak{}\hskip.1em \RB{54.295}{7.752}{2.619}{%
\RG{0.000}{-0.003}{10.9091}{mi}{112}\RG{8.468}{-0.003}{10.9091}{rm}{61}\RG{19.879}{-0.003}{10.9091}{rm}{49}\RG{25.225}{-0.003}{10.9091}{mi}{59}\RG{30.003}{-0.003}{10.9091}{mi}{58}\RG{34.792}{-0.003}{10.9091}{mi}{58}\RG{39.570}{-0.003}{10.9091}{mi}{58}\RG{44.359}{-0.003}{10.9091}{mi}{59}\RG{49.148}{-0.003}{10.9091}{mi}{114}%
}\hskip.1em\allowbreak{} ），其余各行保持不变，所得矩阵记为\allowbreak{}\hskip.1em \RB{8.649}{8.177}{0.732}{%
\RG{0.000}{-0.003}{10.9091}{mi}{67}%
}\hskip.1em\allowbreak{} 。按第\allowbreak{}\hskip.1em \RB{5.758}{7.948}{0.503}{%
\RG{0.000}{-0.003}{10.9091}{mi}{73}%
}\hskip.1em\allowbreak{} 组行作Laplace展开，只有第\allowbreak{}\hskip.1em \RB{7.253}{7.948}{0.732}{%
\RG{0.000}{-0.003}{10.9091}{mi}{74}%
}\hskip.1em\allowbreak{} 组列对应的项可能非零，且\allowbreak{}\hskip.1em \RB{54.619}{8.679}{3.230}{%
\RG{0.000}{-0.003}{10.9091}{mi}{67}\RG{8.425}{-0.003}{10.9091}{rm}{91}\RG{11.456}{-0.003}{10.9091}{mi}{73}\RG{17.020}{-0.003}{10.9091}{mi}{59}\RG{21.809}{-0.003}{10.9091}{mi}{74}\RG{28.725}{-0.003}{10.9091}{rm}{93}\RG{34.714}{-0.003}{10.9091}{rm}{61}\RG{46.060}{-0.003}{10.9091}{mi}{73}\RG{50.772}{-1.639}{7.9701}{mi}{114}%
}\hskip.1em\allowbreak{} ，故

\[
\RB{293.583}{10.436}{3.240}{%
\RG{61.582}{4.502}{7.9701}{mi}{27}\RG{66.241}{4.502}{7.9701}{rm}{40}\RG{69.235}{4.502}{7.9701}{mi}{73}\RG{73.138}{4.502}{7.9701}{rm}{41}\RG{76.131}{4.502}{7.9701}{rm}{43}\RG{82.120}{4.502}{7.9701}{mi}{27}\RG{86.779}{4.502}{7.9701}{rm}{40}\RG{89.773}{4.502}{7.9701}{mi}{74}\RG{94.659}{4.502}{7.9701}{rm}{41}\RG{133.185}{4.502}{7.9701}{mi}{99}\RG{148.671}{4.502}{7.9701}{mi}{99}\RG{0.000}{-0.002}{10.9091}{rm}{100}\RG{5.935}{-0.002}{10.9091}{rm}{101}\RG{10.680}{-0.002}{10.9091}{rm}{116}\RG{16.654}{-0.002}{10.9091}{mi}{67}\RG{28.110}{-0.002}{10.9091}{rm}{61}\RG{39.445}{-0.002}{10.9091}{rm}{40}\RG{43.602}{-0.002}{10.9091}{sy}{0}\RG{52.087}{-0.002}{10.9091}{rm}{49}\RG{57.433}{-0.002}{10.9091}{rm}{41}\RG{99.924}{-0.002}{10.9091}{rm}{100}\RG{105.858}{-0.002}{10.9091}{rm}{101}\RG{110.604}{-0.002}{10.9091}{rm}{116}\RG{116.578}{-0.002}{10.9091}{mi}{65}\RG{124.585}{-0.002}{10.9091}{rm}{91}\RG{127.615}{-0.002}{10.9091}{mi}{73}\RG{136.972}{-0.002}{10.9091}{mi}{59}\RG{141.750}{-0.002}{10.9091}{mi}{74}\RG{152.458}{-0.002}{10.9091}{rm}{93}\RG{155.425}{-0.002}{10.9091}{mi}{58}\RT{279.425}{-0.002}{10.9091}{(}\RT{283.581}{-0.002}{10.9091}{1}\RT{288.926}{-0.002}{10.9091}{)}%
}
\]
由\allowbreak{}\hskip.1em \RB{53.497}{8.293}{2.205}{%
\RG{0.000}{0.004}{10.9091}{mi}{65}\RG{8.007}{0.004}{10.9091}{mi}{65}\RG{16.015}{3.962}{7.9701}{sy}{3}\RG{23.347}{0.004}{10.9091}{rm}{61}\RG{34.685}{0.004}{10.9091}{mi}{68}\RG{43.849}{0.004}{10.9091}{mi}{73}\RG{48.563}{-1.633}{7.9701}{mi}{110}%
}\hskip.1em\allowbreak{} ，矩阵\allowbreak{}\hskip.1em \RB{20.274}{8.293}{0.725}{%
\RG{0.000}{0.004}{10.9091}{mi}{67}\RG{8.422}{0.004}{10.9091}{mi}{65}\RG{16.434}{3.962}{7.9701}{sy}{3}%
}\hskip.1em\allowbreak{} 的第\allowbreak{}\hskip.1em \RB{6.336}{8.075}{0.605}{%
\RG{0.000}{0.004}{10.9091}{mi}{104}%
}\hskip.1em\allowbreak{} 行在\allowbreak{}\hskip.1em \RB{28.891}{8.075}{0.933}{%
\RG{0.000}{0.004}{10.9091}{mi}{104}\RG{9.179}{0.004}{10.9091}{sy}{50}\RG{19.482}{0.004}{10.9091}{mi}{73}\RG{25.051}{3.962}{7.9701}{mi}{99}%
}\hskip.1em\allowbreak{} 时等于\allowbreak{}\hskip.1em \RB{18.818}{8.512}{3.266}{%
\RG{0.000}{0.004}{10.9091}{mi}{68}\RG{9.164}{0.004}{10.9091}{mi}{101}\RG{14.115}{3.962}{7.9701}{sy}{48}\RG{14.115}{-2.694}{7.9701}{mi}{104}%
}\hskip.1em\allowbreak{} ；第\allowbreak{}\hskip.1em \RB{7.911}{7.671}{3.538}{%
\RG{0.000}{0.004}{10.9091}{mi}{105}\RG{3.644}{-1.633}{7.9701}{mi}{112}%
}\hskip.1em\allowbreak{} 行则等于\allowbreak{}\hskip.1em \RB{11.847}{8.293}{0.500}{%
\RG{0.000}{0.004}{10.9091}{mi}{65}\RG{8.007}{3.962}{7.9701}{sy}{3}%
}\hskip.1em\allowbreak{} 的第\allowbreak{}\hskip.1em \RB{8.687}{7.671}{3.538}{%
\RG{0.000}{0.004}{10.9091}{mi}{106}\RG{4.419}{-1.633}{7.9701}{mi}{112}%
}\hskip.1em\allowbreak{} 行。于是

\[
\RB{102.228}{8.685}{3.235}{%
\RG{20.587}{4.507}{7.9701}{sy}{3}\RG{73.340}{4.507}{7.9701}{sy}{3}\RG{0.000}{0.004}{10.9091}{rm}{40}\RG{4.155}{0.004}{10.9091}{mi}{67}\RG{12.576}{0.004}{10.9091}{mi}{65}\RG{24.890}{0.004}{10.9091}{rm}{41}\RG{29.046}{0.004}{10.9091}{rm}{91}\RG{32.074}{0.004}{10.9091}{mi}{73}\RG{37.637}{0.004}{10.9091}{mi}{59}\RG{42.426}{0.004}{10.9091}{mi}{73}\RG{47.996}{0.004}{10.9091}{rm}{93}\RG{53.985}{0.004}{10.9091}{rm}{61}\RG{65.333}{0.004}{10.9091}{mi}{65}\RG{77.643}{0.004}{10.9091}{rm}{91}\RG{80.673}{0.004}{10.9091}{mi}{74}\RG{86.400}{0.004}{10.9091}{mi}{59}\RG{91.178}{0.004}{10.9091}{mi}{73}\RG{96.754}{0.004}{10.9091}{rm}{93}\RG{99.721}{0.004}{10.9091}{mi}{58}%
}
\]
再按\allowbreak{}\hskip.1em \RB{20.273}{8.289}{0.729}{%
\RG{0.000}{-0.000}{10.9091}{mi}{67}\RG{8.422}{-0.000}{10.9091}{mi}{65}\RG{16.433}{3.958}{7.9701}{sy}{3}%
}\hskip.1em\allowbreak{} 的第\allowbreak{}\hskip.1em \RB{9.409}{7.951}{0.500}{%
\RG{0.000}{-0.000}{10.9091}{mi}{73}\RG{5.569}{3.958}{7.9701}{mi}{99}%
}\hskip.1em\allowbreak{} 组行作Laplace展开，唯一可能非零的项选取第\allowbreak{}\hskip.1em \RB{9.408}{7.951}{0.500}{%
\RG{0.000}{-0.000}{10.9091}{mi}{73}\RG{5.568}{3.958}{7.9701}{mi}{99}%
}\hskip.1em\allowbreak{} 组列，其符号为\allowbreak{}\hskip.1em \RB{63.815}{9.892}{3.238}{%
\RG{0.000}{-0.000}{10.9091}{rm}{40}\RG{4.155}{-0.000}{10.9091}{sy}{0}\RG{12.639}{-0.000}{10.9091}{rm}{49}\RG{17.985}{-0.000}{10.9091}{rm}{41}\RG{22.134}{3.958}{7.9701}{rm}{50}\RG{25.983}{3.958}{7.9701}{mi}{27}\RG{30.643}{3.958}{7.9701}{rm}{40}\RG{33.636}{3.958}{7.9701}{mi}{73}\RG{37.540}{6.515}{5.9776}{mi}{99}\RG{41.031}{3.958}{7.9701}{rm}{41}\RG{47.508}{-0.000}{10.9091}{rm}{61}\RG{58.843}{-0.000}{10.9091}{rm}{49}%
}\hskip.1em\allowbreak{} ，因而

\[
\RB{283.406}{9.293}{3.238}{%
\RG{35.423}{4.504}{7.9701}{sy}{3}\RG{67.413}{4.504}{7.9701}{mi}{110}\RG{72.084}{4.504}{7.9701}{sy}{0}\RG{78.072}{4.504}{7.9701}{mi}{114}\RG{108.692}{4.504}{7.9701}{sy}{3}\RG{0.000}{-0.000}{10.9091}{rm}{100}\RG{5.935}{-0.000}{10.9091}{rm}{101}\RG{10.680}{-0.000}{10.9091}{rm}{116}\RG{14.836}{-0.000}{10.9091}{rm}{40}\RG{18.989}{-0.000}{10.9091}{mi}{67}\RG{27.411}{-0.000}{10.9091}{mi}{65}\RG{39.725}{-0.000}{10.9091}{rm}{41}\RG{46.903}{-0.000}{10.9091}{rm}{61}\RG{58.247}{-0.000}{10.9091}{mi}{68}\RG{84.031}{-0.000}{10.9091}{rm}{100}\RG{89.965}{-0.000}{10.9091}{rm}{101}\RG{94.711}{-0.000}{10.9091}{rm}{116}\RG{100.685}{-0.000}{10.9091}{mi}{65}\RG{112.994}{-0.000}{10.9091}{rm}{91}\RG{116.025}{-0.000}{10.9091}{mi}{74}\RG{121.752}{-0.000}{10.9091}{mi}{59}\RG{126.530}{-0.000}{10.9091}{mi}{73}\RG{132.105}{-0.000}{10.9091}{rm}{93}\RG{135.073}{-0.000}{10.9091}{mi}{58}\RT{269.248}{-0.000}{10.9091}{(}\RT{273.405}{-0.000}{10.9091}{2}\RT{278.750}{-0.000}{10.9091}{)}%
}
\]
另一方面，由\allowbreak{}\hskip.1em \RB{66.896}{9.273}{0.613}{%
\RG{0.000}{-0.004}{10.9091}{rm}{100}\RG{5.935}{-0.004}{10.9091}{rm}{101}\RG{10.680}{-0.004}{10.9091}{rm}{116}\RG{16.655}{-0.004}{10.9091}{mi}{65}\RG{24.662}{3.955}{7.9701}{sy}{3}\RG{31.994}{-0.004}{10.9091}{rm}{61}\RG{43.333}{-0.004}{10.9091}{mi}{68}\RG{52.498}{3.955}{7.9701}{mi}{110}\RG{57.169}{3.955}{7.9701}{sy}{0}\RG{63.157}{3.955}{7.9701}{rm}{49}%
}\hskip.1em\allowbreak{} 及行列式的乘法公式，得

\[
\RB{118.269}{9.826}{3.234}{%
\RG{35.423}{4.508}{7.9701}{sy}{3}\RG{100.801}{4.508}{7.9701}{mi}{110}\RG{105.472}{4.508}{7.9701}{sy}{0}\RG{111.459}{4.508}{7.9701}{rm}{49}\RG{0.000}{0.005}{10.9091}{rm}{100}\RG{5.935}{0.005}{10.9091}{rm}{101}\RG{10.680}{0.005}{10.9091}{rm}{116}\RG{14.836}{0.005}{10.9091}{rm}{40}\RG{18.990}{0.005}{10.9091}{mi}{67}\RG{27.412}{0.005}{10.9091}{mi}{65}\RG{39.725}{0.005}{10.9091}{rm}{41}\RG{46.903}{0.005}{10.9091}{rm}{61}\RG{58.248}{0.005}{10.9091}{rm}{40}\RG{62.405}{0.005}{10.9091}{rm}{100}\RG{68.339}{0.005}{10.9091}{rm}{101}\RG{73.085}{0.005}{10.9091}{rm}{116}\RG{79.055}{0.005}{10.9091}{mi}{67}\RG{87.481}{0.005}{10.9091}{rm}{41}\RG{91.635}{0.005}{10.9091}{mi}{68}\RG{115.762}{0.005}{10.9091}{mi}{58}%
}
\]
与(2)比较，约去非零因子\allowbreak{}\hskip.1em \RB{23.672}{7.955}{0.500}{%
\RG{0.000}{0.004}{10.9091}{mi}{68}\RG{9.165}{3.962}{7.9701}{mi}{110}\RG{13.837}{3.962}{7.9701}{sy}{0}\RG{19.825}{3.962}{7.9701}{mi}{114}%
}\hskip.1em\allowbreak{} ，再代入(1)，得到

\[
\RB{208.976}{10.441}{3.236}{%
\RG{24.661}{4.506}{7.9701}{sy}{3}\RG{87.545}{4.506}{7.9701}{mi}{27}\RG{92.204}{4.506}{7.9701}{rm}{40}\RG{95.197}{4.506}{7.9701}{mi}{73}\RG{99.101}{4.506}{7.9701}{rm}{41}\RG{102.093}{4.506}{7.9701}{rm}{43}\RG{108.083}{4.506}{7.9701}{mi}{27}\RG{112.742}{4.506}{7.9701}{rm}{40}\RG{115.735}{4.506}{7.9701}{mi}{74}\RG{120.622}{4.506}{7.9701}{rm}{41}\RG{159.147}{4.506}{7.9701}{mi}{99}\RG{174.634}{4.506}{7.9701}{mi}{99}\RG{192.372}{4.506}{7.9701}{mi}{114}\RG{196.060}{4.506}{7.9701}{sy}{0}\RG{202.047}{4.506}{7.9701}{rm}{49}\RG{0.000}{0.002}{10.9091}{rm}{100}\RG{5.935}{0.002}{10.9091}{rm}{101}\RG{10.680}{0.002}{10.9091}{rm}{116}\RG{16.654}{0.002}{10.9091}{mi}{65}\RG{28.963}{0.002}{10.9091}{rm}{91}\RG{31.994}{0.002}{10.9091}{mi}{74}\RG{37.721}{0.002}{10.9091}{mi}{59}\RG{42.499}{0.002}{10.9091}{mi}{73}\RG{48.075}{0.002}{10.9091}{rm}{93}\RG{54.064}{0.002}{10.9091}{rm}{61}\RG{65.409}{0.002}{10.9091}{rm}{40}\RG{69.565}{0.002}{10.9091}{sy}{0}\RG{78.050}{0.002}{10.9091}{rm}{49}\RG{83.395}{0.002}{10.9091}{rm}{41}\RG{125.886}{0.002}{10.9091}{rm}{100}\RG{131.821}{0.002}{10.9091}{rm}{101}\RG{136.566}{0.002}{10.9091}{rm}{116}\RG{142.540}{0.002}{10.9091}{mi}{65}\RG{150.548}{0.002}{10.9091}{rm}{91}\RG{153.578}{0.002}{10.9091}{mi}{73}\RG{162.935}{0.002}{10.9091}{mi}{59}\RG{167.713}{0.002}{10.9091}{mi}{74}\RG{178.421}{0.002}{10.9091}{rm}{93}\RG{183.206}{0.002}{10.9091}{mi}{68}\RG{206.349}{0.002}{10.9091}{mi}{59}%
}
\]
即所需结论。上述两次展开在\allowbreak{}\hskip.1em \RB{32.090}{7.952}{0.521}{%
\RG{0.000}{0.001}{10.9091}{mi}{73}\RG{5.569}{3.960}{7.9701}{mi}{99}\RG{12.386}{0.001}{10.9091}{rm}{61}\RG{23.725}{0.001}{10.9091}{bb}{63}%
}\hskip.1em\allowbreak{} 时按零阶行列式的约定仍然成立。再设\allowbreak{}\hskip.1em \RB{32.934}{7.948}{0.732}{%
\RG{0.000}{-0.003}{10.9091}{mi}{68}\RG{14.205}{-0.003}{10.9091}{rm}{61}\RG{27.546}{-0.003}{10.9091}{rm}{48}%
}\hskip.1em\allowbreak{} 。若\allowbreak{}\hskip.1em \RB{66.607}{8.679}{3.241}{%
\RG{0.000}{-0.003}{10.9091}{mi}{114}\RG{5.111}{-0.003}{10.9091}{rm}{40}\RG{9.265}{-0.003}{10.9091}{mi}{65}\RG{17.273}{-0.003}{10.9091}{rm}{41}\RG{26.465}{-0.003}{10.9091}{sy}{20}\RG{39.990}{-0.003}{10.9091}{mi}{110}\RG{49.595}{-0.003}{10.9091}{sy}{0}\RG{61.307}{-0.003}{10.9091}{rm}{50}%
}\hskip.1em\allowbreak{} ，则\allowbreak{}\hskip.1em \RB{8.202}{8.286}{0.503}{%
\RG{0.000}{-0.003}{10.9091}{mi}{65}%
}\hskip.1em\allowbreak{} 的所有\allowbreak{}\hskip.1em \RB{26.289}{7.752}{0.612}{%
\RG{0.000}{-0.003}{10.9091}{mi}{110}\RG{9.604}{-0.003}{10.9091}{sy}{0}\RG{21.316}{-0.003}{10.9091}{rm}{49}%
}\hskip.1em\allowbreak{} 阶子式都为零，所以\allowbreak{}\hskip.1em \RB{39.117}{8.286}{0.732}{%
\RG{0.000}{-0.003}{10.9091}{mi}{65}\RG{8.008}{3.956}{7.9701}{sy}{3}\RG{17.349}{-0.003}{10.9091}{rm}{61}\RG{30.697}{-0.003}{10.9091}{mi}{79}%
}\hskip.1em\allowbreak{} 。若\allowbreak{}\hskip.1em \RB{59.721}{8.679}{3.241}{%
\RG{0.000}{-0.003}{10.9091}{mi}{114}\RG{5.111}{-0.003}{10.9091}{rm}{40}\RG{9.265}{-0.003}{10.9091}{mi}{65}\RG{17.272}{-0.003}{10.9091}{rm}{41}\RG{24.450}{-0.003}{10.9091}{rm}{61}\RG{35.795}{-0.003}{10.9091}{mi}{110}\RG{44.217}{-0.003}{10.9091}{sy}{0}\RG{54.748}{-0.003}{10.9091}{rm}{49}%
}\hskip.1em\allowbreak{} ，则由\allowbreak{}\hskip.1em \RB{43.106}{8.286}{0.732}{%
\RG{0.000}{-0.003}{10.9091}{mi}{65}\RG{8.007}{-0.003}{10.9091}{mi}{65}\RG{16.015}{3.956}{7.9701}{sy}{3}\RG{23.348}{-0.003}{10.9091}{rm}{61}\RG{34.686}{-0.003}{10.9091}{mi}{79}%
}\hskip.1em\allowbreak{} ，\allowbreak{}\hskip.1em \RB{11.847}{8.286}{0.503}{%
\RG{0.000}{-0.003}{10.9091}{mi}{65}\RG{8.007}{3.956}{7.9701}{sy}{3}%
}\hskip.1em\allowbreak{} 的列空间包含在一维空间\allowbreak{}\hskip.1em \RB{24.263}{8.286}{0.612}{%
\RG{0.000}{-0.003}{10.9091}{rm}{107}\RG{5.345}{-0.003}{10.9091}{rm}{101}\RG{10.091}{-0.003}{10.9091}{rm}{114}\RG{16.061}{-0.003}{10.9091}{mi}{65}%
}\hskip.1em\allowbreak{} 中，所以\allowbreak{}\hskip.1em \RB{45.246}{8.679}{3.241}{%
\RG{0.000}{-0.003}{10.9091}{mi}{114}\RG{5.111}{-0.003}{10.9091}{rm}{40}\RG{9.265}{-0.003}{10.9091}{mi}{65}\RG{17.272}{3.956}{7.9701}{sy}{3}\RG{21.575}{-0.003}{10.9091}{rm}{41}\RG{28.758}{-0.003}{10.9091}{sy}{20}\RG{40.274}{-0.003}{10.9091}{rm}{49}%
}\hskip.1em\allowbreak{} 。因此无论哪一种情况，\allowbreak{}\hskip.1em \RB{11.847}{8.286}{0.503}{%
\RG{0.000}{-0.003}{10.9091}{mi}{65}\RG{8.007}{3.956}{7.9701}{sy}{3}%
}\hskip.1em\allowbreak{} 的所有\allowbreak{}\hskip.1em \RB{5.147}{5.308}{0.612}{%
\RG{0.000}{-0.003}{10.9091}{mi}{114}%
}\hskip.1em\allowbreak{} （\allowbreak{}\hskip.1em \RB{24.956}{7.752}{1.997}{%
\RG{0.000}{-0.003}{10.9091}{mi}{114}\RG{8.141}{-0.003}{10.9091}{sy}{21}\RG{19.656}{-0.003}{10.9091}{rm}{50}%
}\hskip.1em\allowbreak{} ）阶子式都为零；而右端因含\allowbreak{}\hskip.1em \RB{22.581}{9.275}{0.503}{%
\RG{0.000}{-0.003}{10.9091}{mi}{68}\RG{9.166}{3.956}{7.9701}{mi}{114}\RG{12.854}{3.956}{7.9701}{sy}{0}\RG{18.842}{3.956}{7.9701}{rm}{49}%
}\hskip.1em\allowbreak{} 也为零，等式仍成立。\allowbreak{}\hskip.1em \RB{8.365}{7.992}{0.503}{%
\RG{0.000}{-0.003}{10.9091}{am}{3}%
}\hskip.1em\allowbreak{}

小结：用标准单位行替换指定行，两次Laplace展开便能把互补子式与伴随矩阵子式联系起来。非异情形可约去行列式；奇异情形则利用\allowbreak{}\hskip.1em \RB{46.501}{8.684}{3.236}{%
\RG{0.000}{0.002}{10.9091}{mi}{114}\RG{5.110}{0.002}{10.9091}{rm}{40}\RG{9.265}{0.002}{10.9091}{mi}{65}\RG{17.272}{3.961}{7.9701}{sy}{3}\RG{21.574}{0.002}{10.9091}{rm}{41}\RG{29.385}{0.002}{10.9091}{sy}{20}\RG{41.528}{0.002}{10.9091}{rm}{49}%
}\hskip.1em\allowbreak{} 单独处理，不能直接约去零因子。

% Original page 5.
\section*{2015A07}
设\allowbreak{}\hskip.1em \RB{8.715}{7.954}{0.726}{%
\RG{0.000}{0.003}{10.9091}{mi}{86}%
}\hskip.1em\allowbreak{} 是\allowbreak{}\hskip.1em \RB{33.059}{8.685}{3.235}{%
\RG{0.000}{0.003}{10.9091}{mi}{77}\RG{10.334}{-1.634}{7.9701}{mi}{110}\RG{15.458}{0.003}{10.9091}{rm}{40}\RG{19.612}{0.003}{10.9091}{mi}{75}\RG{29.483}{0.003}{10.9091}{rm}{41}%
}\hskip.1em\allowbreak{} 的子空间，满足\allowbreak{}\hskip.1em \RB{8.715}{7.954}{0.726}{%
\RG{0.000}{0.003}{10.9091}{mi}{86}%
}\hskip.1em\allowbreak{} 中所有的非零矩阵都是非异阵，证明：

\[
\RB{59.165}{8.068}{2.139}{%
\RG{0.000}{-0.003}{10.9091}{rm}{100}\RG{5.935}{-0.003}{10.9091}{rm}{105}\RG{8.902}{-0.003}{10.9091}{rm}{109}\RG{27.130}{-0.003}{10.9091}{mi}{86}\RG{38.762}{-0.003}{10.9091}{sy}{20}\RG{50.276}{-0.003}{10.9091}{mi}{110}\RG{56.658}{-0.003}{10.9091}{mi}{58}\RG{17.855}{-1.639}{7.9701}{mi}{75}%
}
\]
证明：固定一个非零列向量\allowbreak{}\hskip.1em \RB{34.209}{7.953}{0.935}{%
\RG{0.000}{0.002}{10.9091}{mi}{120}\RG{9.100}{0.002}{10.9091}{sy}{50}\RG{19.404}{0.002}{10.9091}{mi}{75}\RG{29.275}{3.961}{7.9701}{mi}{110}%
}\hskip.1em\allowbreak{} ，定义线性映射

\[
\RB{145.021}{8.679}{3.241}{%
\RG{58.850}{4.502}{7.9701}{mi}{110}\RG{0.000}{-0.003}{10.9091}{rm}{8}\RG{10.735}{-0.003}{10.9091}{rm}{58}\RG{16.742}{-0.003}{10.9091}{mi}{86}\RG{28.374}{-0.003}{10.9091}{sy}{0}\RG{35.050}{-0.003}{10.9091}{sy}{33}\RG{48.979}{-0.003}{10.9091}{mi}{75}\RG{63.975}{-0.003}{10.9091}{mi}{59}\RG{90.578}{-0.003}{10.9091}{rm}{8}\RG{98.291}{-0.003}{10.9091}{rm}{40}\RG{102.446}{-0.003}{10.9091}{mi}{84}\RG{110.191}{-0.003}{10.9091}{rm}{41}\RG{117.369}{-0.003}{10.9091}{rm}{61}\RG{128.714}{-0.003}{10.9091}{mi}{84}\RG{136.448}{-0.003}{10.9091}{mi}{120}\RG{142.514}{-0.003}{10.9091}{mi}{58}%
}
\]
若\allowbreak{}\hskip.1em \RB{46.469}{8.684}{3.236}{%
\RG{0.000}{0.002}{10.9091}{rm}{8}\RG{7.713}{0.002}{10.9091}{rm}{40}\RG{11.869}{0.002}{10.9091}{mi}{84}\RG{19.613}{0.002}{10.9091}{rm}{41}\RG{28.264}{0.002}{10.9091}{rm}{61}\RG{41.082}{0.002}{10.9091}{rm}{48}%
}\hskip.1em\allowbreak{} 且\allowbreak{}\hskip.1em \RB{33.480}{8.313}{2.855}{%
\RG{0.000}{0.002}{10.9091}{mi}{84}\RG{12.248}{0.002}{10.9091}{sy}{54}\RG{12.248}{0.002}{10.9091}{rm}{61}\RG{25.060}{0.002}{10.9091}{mi}{79}%
}\hskip.1em\allowbreak{} ，则由假设\allowbreak{}\hskip.1em \RB{8.049}{7.898}{0.500}{%
\RG{0.000}{0.002}{10.9091}{mi}{84}%
}\hskip.1em\allowbreak{} 可逆，故\allowbreak{}\hskip.1em \RB{77.932}{9.279}{0.727}{%
\RG{0.000}{0.002}{10.9091}{mi}{120}\RG{10.574}{0.002}{10.9091}{rm}{61}\RG{23.385}{0.002}{10.9091}{mi}{84}\RG{31.130}{3.961}{7.9701}{sy}{0}\RG{37.117}{3.961}{7.9701}{rm}{49}\RG{41.419}{0.002}{10.9091}{mi}{84}\RG{49.165}{0.002}{10.9091}{mi}{120}\RG{59.737}{0.002}{10.9091}{rm}{61}\RG{72.545}{0.002}{10.9091}{rm}{48}%
}\hskip.1em\allowbreak{} ，与\allowbreak{}\hskip.1em \RB{28.768}{8.313}{2.855}{%
\RG{0.000}{0.002}{10.9091}{mi}{120}\RG{10.574}{0.002}{10.9091}{sy}{54}\RG{10.574}{0.002}{10.9091}{rm}{61}\RG{23.381}{0.002}{10.9091}{rm}{48}%
}\hskip.1em\allowbreak{} 矛盾。因此\allowbreak{}\hskip.1em \RB{57.216}{8.684}{3.225}{%
\RG{0.000}{0.002}{10.9091}{rm}{107}\RG{5.345}{0.002}{10.9091}{rm}{101}\RG{10.091}{0.002}{10.9091}{rm}{114}\RG{16.058}{0.002}{10.9091}{rm}{8}\RG{26.804}{0.002}{10.9091}{rm}{61}\RG{38.144}{0.002}{10.9091}{sy}{102}\RG{43.598}{0.002}{10.9091}{mi}{79}\RG{52.058}{0.002}{10.9091}{sy}{103}%
}\hskip.1em\allowbreak{} ，即\allowbreak{}\hskip.1em \RB{7.613}{7.953}{0.500}{%
\RG{0.000}{0.002}{10.9091}{rm}{8}%
}\hskip.1em\allowbreak{} 为单射。由维数公式，

\[
\RB{303.744}{8.197}{2.139}{%
\RG{151.118}{4.502}{7.9701}{mi}{110}\RG{0.000}{-0.003}{10.9091}{rm}{100}\RG{5.935}{-0.003}{10.9091}{rm}{105}\RG{8.902}{-0.003}{10.9091}{rm}{109}\RG{27.131}{-0.003}{10.9091}{mi}{86}\RG{38.762}{-0.003}{10.9091}{rm}{61}\RG{50.096}{-0.003}{10.9091}{rm}{100}\RG{56.031}{-0.003}{10.9091}{rm}{105}\RG{58.998}{-0.003}{10.9091}{rm}{109}\RG{77.231}{-0.003}{10.9091}{rm}{73}\RG{81.082}{-0.003}{10.9091}{rm}{109}\RG{91.849}{-0.003}{10.9091}{rm}{8}\RG{102.602}{-0.003}{10.9091}{sy}{20}\RG{114.116}{-0.003}{10.9091}{rm}{100}\RG{120.051}{-0.003}{10.9091}{rm}{105}\RG{123.018}{-0.003}{10.9091}{rm}{109}\RG{141.247}{-0.003}{10.9091}{mi}{75}\RG{159.273}{-0.003}{10.9091}{rm}{61}\RG{170.611}{-0.003}{10.9091}{mi}{110}\RG{176.993}{-0.003}{10.9091}{mi}{58}\RG{295.379}{-0.003}{10.9091}{am}{3}\RG{17.855}{-1.639}{7.9701}{mi}{75}\RG{67.955}{-1.639}{7.9701}{mi}{75}\RG{131.973}{-1.639}{7.9701}{mi}{75}%
}
\]
小结：把“每个非零矩阵都可逆”转化为映射\allowbreak{}\hskip.1em \RB{39.283}{7.898}{0.618}{%
\RG{0.000}{0.002}{10.9091}{mi}{84}\RG{11.137}{0.002}{10.9091}{sy}{55}\RG{11.137}{0.002}{10.9091}{sy}{33}\RG{25.441}{0.002}{10.9091}{mi}{84}\RG{33.175}{0.002}{10.9091}{mi}{120}%
}\hskip.1em\allowbreak{} 的单射性，便可将原来的矩阵空间嵌入\allowbreak{}\hskip.1em \RB{6.555}{5.313}{0.607}{%
\RG{0.000}{0.002}{10.9091}{mi}{110}%
}\hskip.1em\allowbreak{} 维列向量空间，从而得到维数上界。


\end{document}
